
\dictentry{Ablative Material}{A material engineered to erode, char, and decompose when exposed to extreme heat, carrying thermal energy away from the underlying structure. Ablative materials are used on spacecraft heat shields, rocket nozzle liners, and re-entry vehicle surfaces where temperatures can exceed thousands of degrees. The material's controlled recession rate is carefully designed so that it sacrifices itself gradually to protect the vehicle during its brief but intense thermal exposure.}{catAE}

\dictentry{Acceptance Test}{A series of standardized tests performed on each production unit to verify that it meets all design specifications and performance requirements before delivery. Acceptance tests typically include functional checks, dimensional inspections, and performance measurements against predefined criteria documented in the type certificate data sheet. A unit that fails acceptance testing is returned for rework or rejection, ensuring that only conforming products reach the customer.}{catAE}

\dictentry{Active Flow Control}{Techniques that use energy to manipulate the airflow over a surface, including synthetic jets, plasma actuators, and oscillating micro-flaps. Active flow control can delay flow separation, reduce drag, and increase lift without the weight penalty of traditional mechanical devices. Unlike passive devices such as vortex generators, active systems can adapt in real time to changing flight conditions, making them promising for next-generation aircraft designs.}{catAE}

\dictentry{Additive Manufacturing}{A fabrication process in which three-dimensional objects are created by depositing material layer by layer from a digital model. In aerospace, additive manufacturing enables the production of complex geometries that are impossible with traditional subtractive methods, reducing weight and material waste. Techniques such as selective laser melting and fused deposition modeling are used to produce engine components, structural brackets, and satellite parts directly from CAD designs.}{catAE}

\dictentry{Aerodynamic Balance}{A design technique in which a portion of a control surface is placed ahead of the hinge line so that aerodynamic forces partially counteract the deflection force required by the pilot. By reducing the hinge moment, aerodynamic balance decreases the control effort needed to move surfaces such as elevators, ailerons, and rudders. This approach is carefully tuned to avoid overbalance, which could cause uncontrollable surface flutter or reversal at high speeds.}{catAE}

\dictentry{Aerospace Engineering}{A branch of engineering concerned with the design, development, testing, and production of aircraft, spacecraft, satellites, and missiles. It encompasses aeronautical engineering for atmospheric flight and astronautical engineering for spaceflight, drawing on disciplines such as aerodynamics, propulsion, structures, and control systems. Aerospace engineers work across the entire vehicle lifecycle from conceptual design through certification, manufacturing, and in-service support.}{catAE}

\dictentry{Afterburner}{A secondary combustion chamber in a jet engine that injects additional fuel directly into the exhaust stream to produce a large temporary increase in thrust. Afterburners are used primarily in military fighter aircraft for takeoff, combat maneuvers, and supersonic flight where peak thrust is needed for short durations. The trade-off is extremely high fuel consumption, so afterburners are typically limited to minutes of continuous operation.}{catAE}

\dictentry{Air Conditioning}{An environmental control subsystem that regulates cabin temperature, humidity, and airflow to maintain passenger and crew comfort during flight. In commercial aircraft, air conditioning is typically supplied by bleed air from the engine compressors, which is cooled through air cycle machines and distributed via an environmental control system. Proper air conditioning is critical not only for comfort but also for preventing fogging of cockpit instruments and protecting sensitive avionics from overheating.}{catAE}

\dictentry{Air Data Computer}{An avionics unit that computes critical flight parameters including indicated airspeed, Mach number, altitude, and total air temperature from sensor inputs such as pitot and static pressures. The air data computer applies corrections for compressibility, instrument errors, and position errors to provide accurate values to the flight instruments and autopilot. Modern air data computers integrate with the flight management system and are essential for safe flight envelope protection.}{catAE}

\dictentry{Aircraft}{Any vehicle designed to travel through the atmosphere using aerodynamic forces generated by wings or other lifting surfaces. Aircraft range from small unmanned drones to large commercial airliners and military fighters, and are classified by their method of lift generation, propulsion type, and intended use. The safe operation of aircraft requires rigorous certification standards covering structural integrity, systems reliability, and flight performance.}{catAE}

\dictentry{Aircraft Wiring}{The network of electrical wires and connectors routing power and signals throughout an aircraft, designed to withstand vibration, temperature extremes, and electromagnetic interference. Aircraft wiring is classified by function into power, signal, and data circuits, each with specific shielding and insulation requirements per standards such as AS50881. Proper wire routing, labeling, and maintenance are critical for preventing electrical faults that could lead to system failures or fires.}{catAE}

\dictentry{Airlock}{A chamber with two sets of sealed doors that allows passage between areas of different atmospheric pressure without equalizing them. In spacecraft, airlocks serve as the transition zone between the pressurized habitable module and the vacuum of space during extravehicular activities. The inner and outer doors are interlocked so that only one can be open at a time, preventing catastrophic depressurization of the crew cabin.}{catAE}

\dictentry{Airplane}{A powered fixed-wing heavier-than-air craft that generates lift through the forward motion of its wings through the atmosphere. Airplanes are the most common form of aircraft, ranging from single-seat light aircraft to wide-body commercial transports carrying hundreds of passengers. They rely on a combination of lift, thrust, weight, and drag forces, controlled by ailerons, elevators, and rudder surfaces.}{catAE}

\dictentry{Airship}{A powered, steerable lighter-than-air craft that uses buoyant gas such as helium to stay aloft and engines for propulsion and directional control. Unlike balloons which drift with the wind, airships can be actively navigated and are used for surveillance, advertising, and tourism applications. Modern airships feature non-rigid or semi-rigid structures and offer long endurance with relatively low fuel consumption compared to heavier-than-air aircraft.}{catAE}

\dictentry{Airworthiness}{The condition of an aircraft conforming to its approved design and being in a state fit for safe operation, as certified by the relevant aviation authority. Airworthiness is established through type certification of the design and maintained through ongoing inspections, maintenance programs, and airworthiness directives issued in response to discovered deficiencies. An aircraft that does not meet airworthiness standards may be grounded until the required corrective actions are completed.}{catAE}

\dictentry{Alternator}{An AC generator driven by the aircraft engine through a belt or gear system that converts mechanical rotation into electrical power for avionics, lighting, and other systems. Alternators are preferred over DC generators in modern aircraft because they are lighter, more reliable, and can produce higher output at lower rotational speeds. The alternator's output is regulated by a voltage regulator and rectified to DC where needed for battery charging and DC-powered equipment.}{catAE}

\dictentry{Aluminum Alloy}{A lightweight metallic alloy with aluminum as the primary element, extensively used in aircraft fuselages, wings, and structural components due to its favorable strength-to-weight ratio and corrosion resistance. Common aerospace grades such as 2024 and 7075 are heat-treated to achieve high tensile strength while remaining formable for sheet metal construction. Aluminum alloys are being gradually supplemented by composite materials in modern aircraft, but they remain the dominant structural material due to their predictable fatigue behavior and repairability.}{catAE}

\dictentry{Anti-Servo Tab}{A tab attached to a control surface that moves in the same direction as the surface deflection, increasing the hinge moment to provide artificial feel and prevent over-controlling. The anti-servo tab adds resistance proportional to the control input, giving the pilot tactile feedback about the magnitude of the command being applied. This is particularly important in aircraft with aerodynamically balanced control surfaces where natural feedback is reduced.}{catAE}

\dictentry{Anti-Skid System}{A safety system that prevents wheel lockup during braking by automatically modulating brake pressure, analogous to anti-lock brakes on automobiles. The system uses wheel speed sensors to detect imminent lockup and rapidly pulses the brakes to maintain optimal tire-to-runway friction. Anti-skid systems significantly reduce landing distance and prevent tire blowouts, especially on wet or contaminated runways.}{catAE}

\dictentry{APU}{Auxiliary Power Unit, a small gas turbine engine typically housed in the tail section that provides compressed air for engine starting and electrical power while the main engines are off. The APU allows the aircraft to operate self-sufficiently on the ground without relying on external ground power units or air start carts. It is also used in flight as a backup power source and to maintain cabin pressurization in the event of an engine failure.}{catAE}

\dictentry{ARINC 429}{A widely adopted standard defining the data bus communication protocol used in commercial avionics systems for transmitting digital information between avionics components. The standard specifies a twisted-pair wire configuration, a 32-bit word format with label, data, and parity fields, and a simplex broadcast architecture where one transmitter sends to multiple receivers. ARINC 429 has been the backbone of commercial avionics data communication since its introduction in the 1980s.}{catAE}

\dictentry{Attitude Control}{The management of a spacecraft's or aircraft's orientation in three-dimensional space using actuators such as reaction wheels, control moment gyros, thrusters, or aerodynamic surfaces. Proper attitude control is essential for pointing antennas toward Earth, orienting solar panels toward the Sun, and aligning sensors toward their targets. The control system continuously compares the current orientation to the desired attitude and generates corrective torques to minimize pointing errors.}{catAE}

\dictentry{Attitude Determination}{The process of calculating a spacecraft's current orientation in space by processing measurements from sensors such as star trackers, sun sensors, magnetometers, and inertial measurement units. Attitude determination algorithms, including TRIAD, QUEST, and extended Kalman filters, fuse multi-source sensor data to produce accurate three-axis attitude estimates. Accurate attitude determination is a prerequisite for attitude control, navigation, and payload pointing operations.}{catAE}

\dictentry{Autoclave Curing}{A composite manufacturing process in which fiber-reinforced polymer parts are cured under elevated temperature and pressure inside a sealed vessel called an autoclave. The combination of heat and external pressure ensures thorough consolidation of the laminate layers, minimal void content, and optimal fiber-to-resin ratio. Autoclave-cured composites achieve the highest quality and are specified for primary aircraft structures such as fuselage panels and wing skins.}{catAE}

\dictentry{Autonomous Flight}{Flight operations conducted without direct human intervention, relying on onboard computers and sensors to plan, navigate, and control the vehicle's trajectory. Autonomous flight systems use artificial intelligence, GPS, lidar, and computer vision to perceive the environment and make real-time decisions. While fully autonomous passenger aircraft are still in development, autonomous capabilities are already standard in military drones and increasingly used in general aviation for autopilot and collision avoidance.}{catAE}

\dictentry{Autopilot}{An automated flight control system that maintains the aircraft's flight path, altitude, heading, and airspeed without continuous pilot input. Modern autopilots interface with the flight management system, air data computers, and inertial reference systems to execute complex navigational procedures including instrument approaches. The autopilot reduces pilot workload during long cruise segments but requires constant monitoring and can be disconnected by the pilot at any time.}{catAE}

\dictentry{Autorotation Landing}{A helicopter landing technique performed after engine failure in which the rotor is driven entirely by upward airflow through the blade disc as the aircraft descends. The pilot manages rotor RPM by controlling collective pitch and descent rate, then flares just before touchdown to convert descent energy into a cushion of air for a soft landing. Autorotation is a fundamental skill that all helicopter pilots must demonstrate during certification.}{catAE}

\dictentry{Autothrottle}{An automatic system that adjusts engine thrust settings to maintain a target airspeed, Mach number, or vertical speed without manual throttle inputs from the pilot. The autothrottle uses inputs from the air data computer and flight management system to continuously modulate thrust lever position. It is especially useful during takeoff, climb, cruise, and approach phases, reducing pilot workload and improving fuel efficiency through precise thrust management.}{catAE}

\dictentry{Balance Tab}{A tab attached to a control surface that deflects in the opposite direction to the surface, creating an aerodynamic force that helps reduce the hinge moment. The balance tab assists the pilot by partially neutralizing the aerodynamic resistance of the control surface, making it easier to deflect. It is mechanically linked to the fixed structure so that its movement is always opposite to that of the control surface.}{catAE}

\dictentry{Balloon}{An unpowered lighter-than-air craft that rises and floats using buoyant gas such as helium or hot air, with altitude controlled by releasing gas or dropping ballast. Unlike airships, balloons have no propulsion or steering capability and drift with ambient wind currents. They are used for scientific research, weather observation, and recreational purposes, with some reaching the stratosphere at altitudes exceeding 30 kilometers.}{catAE}

\dictentry{Battery}{An electrochemical power storage device that converts chemical energy into electrical energy through reversible reactions, used in aircraft for starting, emergency power, and in electric aircraft for primary propulsion. Aerospace batteries must meet stringent requirements for weight, reliability, and performance across wide temperature and pressure ranges. Lithium-ion batteries have largely replaced older nickel-cadmium types due to their higher energy density, though thermal management remains a critical safety concern.}{catAE}

\dictentry{Bellcrank}{A lever mechanism in older aircraft that changes the direction of control cable movement from the cockpit to the control surfaces. This L-shaped or triangular component converts a pulling motion in one direction to a pulling motion at a different angle, enabling routing of control cables around structural obstacles. Bellcranks remain in use on many light and vintage aircraft due to their simplicity and mechanical reliability.}{catAE}

\dictentry{Bleed Air}{Hot, compressed air tapped from the compressor stages of a gas turbine engine and routed through the pneumatic system for cabin pressurization, air conditioning, wing and engine anti-icing, and engine starting. Bleed air extraction reduces engine thrust output and increases specific fuel consumption, which is why some modern designs like the Boeing 787 use electrically driven compressors instead. The bleed air must pass through pre-coolers and valves to regulate its temperature and pressure before reaching downstream systems.}{catAE}

\dictentry{Blown Flap}{A high-lift device that augments the aerodynamic effect of a flap by ejecting high-velocity engine bleed air or fan air over its upper surface, energizing the boundary layer and delaying flow separation. Blown flaps can significantly increase the maximum lift coefficient, allowing aircraft to take off and land at lower speeds and shorter runway distances. This technology was employed on aircraft such as the F-4 Phantom and the XC-142 tilt-wing STOL transport.}{catAE}

\dictentry{Boosted Controls}{Flight control systems that use hydraulic, pneumatic, or electric actuators to amplify pilot inputs and reduce the physical effort required to move control surfaces. Boosted controls are essential on large and high-speed aircraft where aerodynamic loads on control surfaces would be too great for unassisted human strength. They typically include artificial feel systems that restore tactile feedback so the pilot retains a sense of the aerodynamic forces acting on the controls.}{catAE}

\dictentry{Booster}{The initial propulsion stage of a launch vehicle, providing the high thrust needed to lift the rocket off the launch pad and through the dense lower atmosphere. Boosters may be solid or liquid fueled and are often jettisoned once their propellant is exhausted to reduce the remaining vehicle mass. Strap-on boosters supplement the core stage thrust and are commonly used on vehicles such as the Atlas V, Delta IV, and Ariane 5.}{catAE}

\dictentry{Brake}{A mechanical device on the aircraft landing gear that converts kinetic energy into heat through friction to decelerate the aircraft during taxi, landing rollout, and rejected takeoff. Aircraft brakes typically use carbon or steel disc stacks clamped by hydraulic pistons, with multi-disc configurations providing sufficient stopping power for high-energy scenarios. Many modern brakes incorporate anti-skid systems and are rated for energy absorption in terms of kinetic energy at a specified touchdown speed.}{catAE}

\dictentry{Braking System}{The complete assembly of components responsible for decelerating the aircraft on the runway, including wheel brakes, anti-skid controllers, brake-by-wire electronics, and the autobrake system. The braking system must dissipate enormous amounts of kinetic energy, often exceeding 30 megajoules for a wide-body aircraft landing at maximum weight. Redundant hydraulic or electric actuation ensures braking capability is maintained even in the event of a primary system failure.}{catAE}

\dictentry{Bulkhead}{A transverse structural partition within the fuselage that carries concentrated loads from engines, landing gear, and control surface attachments, and also serves as a pressure seal between cabin zones. Bulkheads are typically made from thick aluminum plate or composite layups and are reinforced at attachment points with local doubler plates. In pressurized aircraft, the forward and aft pressure bulkheads are critical structural elements that contain the cabin differential pressure.}{catAE}

\dictentry{Bus Bar}{A conductive metal strip or bar used in aircraft electrical systems to distribute power from generators and batteries to multiple circuits through a centralized connection point. Bus bars are typically housed in electrical load management centers and are sized to carry the maximum anticipated current while minimizing voltage drop. They provide a simpler and more reliable alternative to individual wire runs for high-current power distribution throughout the aircraft.}{catAE}

\dictentry{Cabin Pressurization}{The process of maintaining a breathable, comfortable atmospheric pressure inside the aircraft cabin during flight at high altitudes where ambient air pressure is too low for human survival. A pressurization system compresses outside air and controls its release through outflow valves to maintain a cabin altitude typically equivalent to 6,000 to 8,000 feet. Structural design must account for the pressure differential between the cabin and the outside atmosphere, which places cyclic loads on the fuselage skin and frames.}{catAE}

\dictentry{Cable}{A wire rope used in aircraft to transmit mechanical control inputs from the cockpit to the control surfaces through a system of pulleys and guides. Steel cables are favored for their high tensile strength, flexibility, and resistance to fatigue from repeated bending around pulleys. Cable systems require regular inspection for fraying, corrosion, and proper tension, as a cable failure can result in loss of control of the affected surface.}{catAE}

\dictentry{Calibrated Airspeed}{The indicated airspeed corrected for instrument errors and position errors caused by the placement of the pitot tube and static ports on the aircraft. Calibrated airspeed is used as the primary reference for aircraft performance calculations including takeoff and landing speeds, climb rates, and stall margins. It differs from true airspeed by the density ratio and is the speed the pilot sees after applying known corrections to the raw instrument reading.}{catAE}

\dictentry{Canopy}{The transparent enclosure over the cockpit, typically made from stretched acrylic or polycarbonate, that provides aerodynamic streamlining while protecting the crew from wind blast, noise, and environmental hazards. Canopies are designed to be jettisoned or ejected in emergency situations, and many feature built-in detonation cord systems for fragmentation before ejection. The optical quality of the canopy material is critical, as distortion can impair the pilot's ability to acquire visual targets or judge distances during landing.}{catAE}

\dictentry{Carbon Fiber}{A strong, stiff, lightweight composite material made from thin filaments of carbon atoms bonded in a crystalline structure, used extensively in aerospace for structural components. Carbon fiber reinforced polymer offers a strength-to-weight ratio significantly exceeding that of aluminum and can be oriented in specific directions to optimize load-bearing capability. It is used in fuselage sections, wing skins, control surfaces, and interior structures, though it requires specialized repair techniques and conducts electricity less predictably than metals.}{catAE}

\dictentry{Cat A Takeoff}{A critical engine failure speed defined in certification regulations where the aircraft must be able to continue a safe takeoff even if the most critical engine fails at or after reaching this speed. If the engine fails below V1 the takeoff must be rejected, but at or above V1 the pilot is committed to flying and must achieve a safe climb gradient with the remaining engines. Cat A establishes the minimum climb performance required for one-engine-inoperative scenarios during takeoff.}{catAE}

\dictentry{Cat B Takeoff}{A takeoff classification in which the aircraft is not required to demonstrate one-engine-inoperative climb performance after V1, because it has already demonstrated sufficient acceleration and climb capability on all engines. Cat B applies to lighter aircraft or specific operational conditions where the takeoff safety margins are met without the need for OEI procedures. This category simplifies takeoff performance calculations and allows higher obstacle clearance speeds.}{catAE}

\dictentry{Caution Warning System}{An avionics subsystem that monitors aircraft systems and alerts the flight crew to abnormal or potentially hazardous conditions through visual annunciations and aural warnings. The system categorizes alerts into warning (red), caution (amber), and advisory (informational) levels, each requiring an appropriate crew response. Modern caution warning systems are integrated with electronic centralised aircraft monitoring to provide system synoptic displays and help crews quickly identify and respond to malfunctions.}{catAE}

\dictentry{Ceiling}{The maximum density altitude at which an aircraft can maintain a specified rate of climb, typically 100 feet per minute for multi-engine aircraft with one engine inoperative. The service ceiling represents the practical upper limit of the aircraft's operational altitude where climb performance degrades to near zero. Ceiling is affected by aircraft weight, atmospheric temperature, and engine performance, and it varies significantly between day and night conditions.}{catAE}

\dictentry{Ceramic Matrix Composite}{A material consisting of ceramic fibers embedded in a ceramic matrix, designed to withstand extremely high temperatures exceeding 1,200 degrees Celsius where metal alloys and polymer composites would fail. CMCs offer low density, high strength retention at temperature, and resistance to oxidation, making them ideal for turbine engine hot-section components such as shrouds, vanes, and nozzle flaps. The use of CMCs in jet engines reduces cooling air requirements and improves thermodynamic efficiency.}{catAE}

\dictentry{Circuit Breaker}{A resettable overcurrent protection device that automatically opens the circuit when excessive current is detected, preventing wire overheating and potential fires. Unlike a fuse, which must be replaced after blowing, a circuit breaker can be manually reset once the fault condition is cleared. In aircraft, circuit breakers are grouped on panels accessible to the crew and are labeled to correspond with specific electrical loads for easy identification during troubleshooting.}{catAE}

\dictentry{Circular Orbit}{An orbit in which the satellite maintains a constant altitude above the celestial body, resulting in an orbital path that traces a perfect circle with an eccentricity of exactly zero. The orbital velocity in a circular orbit is constant and is determined solely by the gravitational parameter of the central body and the orbital radius. Many operational satellites including communications and weather satellites use circular orbits to maintain consistent coverage and predictable ground tracks.}{catAE}

\dictentry{Circulation Control Wing}{A concept where high-velocity air is blown over a rounded trailing edge of the wing to control lift and reduce drag by manipulating the circulation around the airfoil. By adjusting the blowing rate, pilots can achieve large changes in lift coefficient without mechanically moving control surfaces. The Coanda effect causes the blown jet to follow the curved surface, effectively changing the wing's aerodynamic shape in flight.}{catAE}

\dictentry{CNC Machining}{Computer numerical controlled machining, a manufacturing process in which pre-programmed computer software controls the movement of machine tools to produce precision parts from metal or composite stock. CNC mills, lathes, and routers can achieve tolerances of thousandths of an inch and produce complex geometries directly from CAD models. In aerospace, CNC machining is essential for producing engine components, structural fittings, and prototype parts with high repeatability and minimal human error.}{catAE}

\dictentry{Cockpit}{The compartment at the front of an aircraft from which the pilot controls the flight, housing the flight instruments, control columns, throttle quadrant, and avionics displays. In modern glass cockpits, traditional analog gauges have been replaced by large multi-function electronic displays that present flight, navigation, and systems data. The cockpit layout is designed for ergonomic efficiency and situational awareness, with critical instruments positioned within the pilot's primary field of view.}{catAE}

\dictentry{Collective}{The flight control in a helicopter that simultaneously changes the pitch angle of all main rotor blades by the same amount, controlling the total lift generated by the rotor. Raising the collective increases blade pitch and generates more lift, while lowering it reduces lift and initiates descent. The collective is linked to the engine throttle or governor to maintain rotor RPM as power demand changes with pitch setting.}{catAE}

\dictentry{Collective Pitch}{A helicopter flight control condition in which the pitch angle of every main rotor blade is changed simultaneously and by the same amount, directly regulating the total aerodynamic lift produced by the rotor disc. Increasing collective pitch increases blade angle of attack, producing greater lift but also demanding more engine power to maintain rotor speed. The collective pitch lever in the cockpit is the primary means of controlling the helicopter's rate of climb or descent.}{catAE}

\dictentry{Combustion Chamber}{The engine section where fuel is injected into the compressed air stream and ignited to produce high-temperature, high-pressure gas that drives the turbine and produces thrust. Combustion chambers are designed to maintain stable and efficient burning across a wide range of altitudes, temperatures, and power settings. Modern annular and can-annular designs achieve combustion efficiencies above 99 percent while minimizing pollutant emissions such as nitrogen oxides and unburned hydrocarbons.}{catAE}

\dictentry{Composite Material}{A material combining two or more constituent materials with significantly different physical properties to produce a material with characteristics superior to either component alone. In aerospace, the most common composites are fiber-reinforced polymers where high-strength fibers such as carbon or glass are embedded in a polymer resin matrix. Composites offer design flexibility, corrosion resistance, and significant weight savings, with modern aircraft like the Boeing 787 and Airbus A350 using over 50 percent composite structures by weight.}{catAE}

\dictentry{Control Column}{The primary flight control input device in the cockpit, typically a column or stick that the pilot moves to command pitch and roll. Forward and backward movement of the column commands nose-down and nose-up pitch through elevator or stabilator deflection, while lateral movement commands roll through aileron or spoiler deflection. Modern fly-by-wire aircraft may use sidesticks or center sticks that send electrical signals rather than mechanical linkages to the control surfaces.}{catAE}

\dictentry{Control Horn}{A rigid arm or bracket attached to a control surface that provides an attachment point for pushrods, cables, or linkages that transmit control forces. The geometry of the control horn determines the mechanical advantage and deflection ratio between the pilot's input and the resulting surface movement. Proper sizing and positioning of control horns is critical for achieving the desired control authority and response characteristics.}{catAE}

\dictentry{Control Law}{A mathematical algorithm or set of equations implemented in a fly-by-wire flight control system that determines how pilot inputs and sensor data are translated into commands sent to the control surface actuators. Control laws define the relationship between cockpit input and aircraft response, incorporating gain scheduling, load factor limiting, angle of attack protection, and stability augmentation. Different control law modes such as normal, direct, and degraded are available depending on the aircraft's configuration and system status.}{catAE}

\dictentry{Control Surface}{A movable aerodynamic surface on an aircraft used to control the vehicle's rotation about its center of gravity in pitch, roll, and yaw. Primary control surfaces include elevators for pitch, ailerons for roll, and the rudder for yaw, while secondary surfaces such as trim tabs and spoilers provide supplementary control. Control surfaces are deflected by the pilot through mechanical linkages, hydraulic actuators, or electrical signals in fly-by-wire systems.}{catAE}

\dictentry{Cosmic Ray}{High-energy charged particles, primarily protons and atomic nuclei, that travel through space at nearly the speed of light and pose radiation hazards to spacecraft electronics and crew. Galactic cosmic rays originate from supernovae and other energetic astrophysical events, while solar cosmic rays are emitted during solar flares and coronal mass ejections. Spacecraft shielding, radiation-hardened electronics, and mission timing strategies are used to mitigate the effects of cosmic ray exposure on long-duration missions.}{catAE}

\dictentry{Cyclic}{The control stick in a helicopter that tilts the rotor disc in the desired direction of flight by varying blade pitch cyclically during each rotation. Moving the cyclic forward tilts the rotor disc forward, causing the helicopter to accelerate in that direction. The cyclic works through a swashplate mechanism that translates the pilot's stick inputs into sinusoidal pitch variations around the rotor azimuth.}{catAE}

\dictentry{Cyclic Pitch}{The cyclically varying pitch angle of helicopter rotor blades that changes once per revolution, creating an asymmetric lift distribution that tilts the rotor disc in the desired direction of flight. The blade at a given azimuth position has a specific pitch that produces more lift on one side of the rotor than the other, generating a net horizontal force. Cyclic pitch is controlled by the pilot's cyclic stick through a swashplate assembly that converts linear stick motion into the required sinusoidal blade pitch variation.}{catAE}

\dictentry{Data Bus}{A shared communication network within an aircraft that allows multiple avionics systems and Line Replaceable Units to exchange digital data using a standardized protocol. Common data bus standards include ARINC 429 for commercial aircraft and MIL-STD-1553 for military applications, each defining electrical characteristics, word formats, and transmission rates. The data bus architecture reduces wiring weight and complexity by replacing dedicated point-to-point wiring with a networked approach.}{catAE}

\dictentry{Dead Reckoning}{A navigation method that estimates the current position of an aircraft by advancing a known previous position along a course at the groundspeed for the elapsed time. Dead reckoning relies on accurate heading, airspeed, and wind information, and its accuracy degrades over time due to cumulative errors from wind estimation, instrument inaccuracies, and heading drift. Despite the availability of GPS and inertial navigation, dead reckoning remains a fundamental backup technique taught to all pilots.}{catAE}

\dictentry{Delta-v}{The total change in velocity that a spacecraft can achieve through its propulsion system, expressed in meters per second and representing the fundamental measure of a spacecraft's maneuvering capability. Delta-v is determined by the Tsiolkovsky rocket equation, which relates it to the specific impulse of the propulsion system and the ratio of initial to final mass. Mission planners allocate delta-v budgets across phases such as orbit insertion, station-keeping, and deorbit to ensure sufficient propellant is carried.}{catAE}

\dictentry{Digital Twin}{A virtual replica of a physical aircraft or system that is continuously updated with real-time sensor data and maintained throughout the product lifecycle for simulation, prediction, and optimization. Digital twins enable engineers to test design changes, predict maintenance needs, and diagnose faults without modifying the actual hardware. In aerospace, digital twins are used for structural health monitoring, engine performance optimization, and fleet-wide analytics to reduce downtime and improve safety.}{catAE}

\dictentry{Dissymmetry of Lift}{An aerodynamic condition in helicopter forward flight where the advancing rotor blade experiences higher relative airflow and generates more lift than the retreating blade, which operates at a lower effective airspeed. This lift imbalance creates a rolling tendency that the helicopter must counteract through blade flapping, which automatically adjusts blade pitch cyclically to equalize lift across the rotor disc. Dissymmetry of lift is a fundamental constraint on helicopter forward speed, as the retreating blade eventually stalls at high airspeeds.}{catAE}

\dictentry{DO-178C}{Software Considerations in Airborne Systems and Equipment Certification, the international standard that defines the processes, activities, and documentation required for developing airborne software with the appropriate level of rigor for its criticality. The standard assigns software levels from DAL A (catastrophic failure condition) to DAL E (no effect), with each level requiring progressively more stringent verification, traceability, and configuration management. DO-178C is mandated by aviation authorities worldwide and must be satisfied before any new or modified software can be certified for use in civil aircraft.}{catAE}

\dictentry{DO-254}{Design Assurance Guidance for Airborne Electronic Hardware, the standard that defines the hardware development and verification processes required for airborne electronic hardware to be certified in civil aircraft. DO-254 establishes assurance levels from Level A (catastrophic failure condition) to Level E (no effect), each requiring specific levels of design review, verification, and process rigor. The standard covers custom hardware such as ASICs, FPGAs, and programmable logic devices used in avionics systems.}{catAE}

\dictentry{Docking Port}{A standardized mechanical and electrical interface on a spacecraft designed to allow two vehicles to physically connect, transfer crew, and exchange power, data, and consumables. International docking standards such as the International Docking Adapter ensure interoperability between spacecraft from different nations and programs. Docking ports include capture mechanisms, alignment guides, and sealed pressure vestibules that allow safe crew transfer between pressurized modules.}{catAE}

\dictentry{Downwash}{The component of airflow induced downward by a wing or rotor as it generates lift, which alters the velocity field around and behind the lifting surface. Downwash reduces the effective angle of attack at the tail, influencing longitudinal stability, and affects the performance of following aircraft during formation flight or in airport approach paths. In helicopters, downwash is particularly intense and contributes to phenomena such as ground effect and vortex ring state during descent.}{catAE}

\dictentry{EASA}{The European Union Aviation Safety Agency, responsible for civil aviation safety regulation, type certification, and oversight of airlines and maintenance organizations across EU member states. EASA was established in 2002 and has authority over aircraft design approval, aircrew licensing, and operational standards for European civil aviation. The agency also cooperates with the FAA and other international authorities through bilateral agreements to harmonize certification standards.}{catAE}

\dictentry{ECAM}{Electronic Centralized Aircraft Monitor, an Airbus avionics system that continuously monitors aircraft systems and automatically displays relevant information and procedural guidance to the flight crew in the event of a malfunction. ECAM consolidates engine, warning, and system synoptic displays into a single format, reducing pilot workload during abnormal situations by presenting only the information relevant to the current fault. The system is divided into upper and lower displays that show engine parameters and memo information respectively.}{catAE}


\dictentry{Effective Translational Lift}{The increase in helicopter rotor lift that occurs during forward flight as the rotor moves out of its own downwash and into undisturbed air. At hover, the rotor operates in its own recirculating wake, but as forward speed increases, more of the rotor disc encounters clean airflow, improving aerodynamic efficiency. This effect reduces the power required for a given lift as speed increases from hover, contributing to the characteristic U-shaped power curve of helicopters.}{catAE}

\dictentry{EICAS}{Engine Indicating and Crew Alerting System, a Boeing avionics system that displays engine performance parameters and alerts the flight crew to system malfunctions through visual and aural warnings. EICAS presents critical engine data such as thrust settings, temperatures, and fuel flow on dedicated displays while automatically prioritizing alerts to focus crew attention on the most urgent issues. The system reduces pilot workload during abnormal conditions by presenting information in a structured and prioritized format.}{catAE}

\dictentry{Electric Aircraft}{An aircraft that uses electric motors powered by batteries or fuel cells as its primary propulsion source, potentially offering lower emissions, reduced noise, and simpler maintenance compared to conventional piston or turbine engines. While fully electric aircraft are currently limited to small urban air mobility vehicles and training aircraft due to energy density constraints of batteries, hybrid-electric designs extend the range of viable electric flight. Advances in battery technology, power electronics, and motor design are rapidly expanding the practical boundaries of electric aviation.}{catAE}

\dictentry{Electrical System}{The integrated network of generators, batteries, transformers, wiring, and distribution devices that generates, conditions, and delivers electrical power to all aircraft systems and loads. Aircraft electrical systems typically operate on a combination of AC at 115 volts 400 hertz and DC at 28 volts, with dedicated buses for essential, non-essential, and emergency loads. The system includes multiple layers of protection including circuit breakers, overcurrent relays, and load shedding logic to maintain critical systems during faults.}{catAE}

\dictentry{Electron Beam Welding}{A high-energy welding process that uses a focused beam of high-velocity electrons in a vacuum to join metals by localized melting and solidification. The concentrated energy density produces deep, narrow welds with minimal heat-affected zones, making it ideal for joining thick sections and dissimilar metals in aerospace applications. Electron beam welding is used for turbine components, landing gear parts, and other critical structural joints where high strength and precision are required.}{catAE}

\dictentry{Elliptical Orbit}{An orbit with an eccentricity between zero and one, in which the satellite's altitude varies between a minimum at periapsis and a maximum at apoapsis during each revolution. Elliptical orbits arise naturally during orbital transfers such as the Hohmann transfer and are exploited for specific mission profiles including Molniya communications orbits with high-latitude coverage. The satellite's velocity varies continuously along an elliptical path, being fastest at periapsis and slowest at apoapsis per Kepler's second law.}{catAE}

\dictentry{Empennage}{The tail assembly of an aircraft, consisting of the horizontal stabilizer, vertical stabilizer, and associated control surfaces, which provides longitudinal, directional, and yaw stability. The empennage counteracts the natural tendency of the aircraft to pitch, yaw, or diverge from its flight path by producing corrective aerodynamic forces. In conventional designs the horizontal stabilizer provides a downward force to balance the nose-heavy center of gravity, while the vertical fin provides directional weathercock stability.}{catAE}

\dictentry{Engine Display}{A cockpit display screen that presents real-time engine performance parameters to the flight crew, including thrust setting, exhaust gas temperature, oil pressure, fuel flow, and vibration levels. Engine displays are integrated into the glass cockpit layout and are managed by the EICAS or ECAM system to alert crews to abnormal operating conditions. Modern engine displays use color-coded scales and trend indicators to help pilots quickly assess engine health and detect developing faults before they become critical.}{catAE}

\dictentry{Environmental Control System}{The integrated system that manages cabin pressure, temperature, air quality, and ventilation within an aircraft to maintain a habitable environment for passengers and crew at all altitudes. The ECS uses engine bleed air or electric compressors to provide conditioned air through a network of ducts, valves, and heat exchangers. Beyond comfort, the system must prevent fogging, maintain adequate oxygen levels, and protect avionics bays from overheating during flight.}{catAE}


\dictentry{FAA}{The Federal Aviation Administration, the United States government agency responsible for regulating all aspects of civil aviation including airworthiness certification, pilot licensing, air traffic control, and airport safety standards. The FAA issues type certificates, production certificates, and airworthiness certificates that authorize aircraft to operate in the National Airspace System. It also investigates aviation accidents and incidents through the National Transportation Safety Board and issues airworthiness directives to address safety concerns.}{catAE}

\dictentry{Fail-Safe Design}{A design philosophy in which the structure and systems are engineered so that if a component fails, the remaining structure or redundant systems can carry the loads or maintain functionality without catastrophic consequences. In aircraft structures, fail-safe design relies on multiple load paths, crack arrest features, and safety pins so that failure of one member does not lead to overall structural collapse. This approach contrasts with the older safe-life approach and is a cornerstone of modern airworthiness standards.}{catAE}

\dictentry{Fairing}{An aerodynamic cover or streamlined enclosure that reduces drag by smoothing the airflow over protrusions such as antennas, landing gear, and payload interfaces on a launch vehicle or aircraft. On launch vehicles, the payload fairing protects the satellite or spacecraft from aerodynamic loads, heating, and acoustic vibration during ascent and is jettisoned once the vehicle exits the atmosphere. The fairing halves are typically separated by pyrotechnic or pneumatic mechanisms and fall away from the vehicle on a predictable trajectory.}{catAE}

\dictentry{Fatigue Test}{A structural test performed on aircraft components or full-scale airframes by applying repeated cyclic loads that simulate the stress spectrum experienced during the aircraft's operational life. Fatigue testing is essential for determining crack initiation sites, propagation rates, and the remaining useful life of structural elements. Results from fatigue tests validate analytical fatigue life predictions and establish inspection intervals and maintenance requirements that ensure continued airworthiness throughout the fleet's service life.}{catAE}

\dictentry{Fault Tolerance}{The ability of a system to continue operating correctly, possibly at a reduced performance level, even when one or more of its components have failed. In aerospace, fault tolerance is achieved through redundancy such as dual or triple-redundant flight computers, multiple hydraulic systems, and parallel electrical generators. Certification requirements for transport category aircraft typically mandate that no single failure or combination of failures can result in a catastrophic condition.}{catAE}

\dictentry{Feathering Hinge}{A hinge mechanism on a helicopter rotor blade that allows the blade to rotate about its longitudinal axis to change its pitch angle during flight. The feathering hinge is actuated by the swashplate through pitch links, enabling both collective and cyclic pitch control of the blade. In some rotor head designs, the feathering hinge is combined with flapping and lead-lag hinges to form a fully articulated rotor system.}{catAE}

\dictentry{Fenestron}{A ducted tail rotor design, also known as a fan-in-fin, in which multiple small blades are enclosed within a circular shroud built into the vertical stabilizer. The Fenestron provides anti-torque and directional control like a conventional tail rotor but with reduced noise, lower vulnerability to foreign object damage, and improved ground safety due to the guarded blade tips. Airbus Helicopters has widely adopted the Fenestron on models including the H135, H145, and H160.}{catAE}

\dictentry{Fiberglass}{A composite material made of fine glass fibers embedded in a polymer resin matrix, offering good strength-to-weight ratio, electrical insulation, and corrosion resistance. Fiberglass is widely used in aircraft radomes, fairings, interior panels, and helicopter fuselage shells where its ability to be molded into complex shapes and its transparency to radar signals make it particularly suitable. It is less stiff and strong than carbon fiber but significantly less expensive, making it the material of choice for non-primary structural applications.}{catAE}

\dictentry{Filament Winding}{A composite manufacturing process in which resin-impregnated fibers are precisely wound onto a rotating mandrel in predetermined patterns to create hollow cylindrical or spherical structures. The winding angle and fiber tension are controlled to optimize the structural properties for the intended loading conditions, typically internal pressure or axial loads. Filament winding is used to produce rocket motor casings, pressure vessels, helicopter tail booms, and drive shafts with excellent hoop and axial strength.}{catAE}

\dictentry{Fin}{Another term for the vertical stabilizer, the fixed vertical tail surface of an aircraft that provides directional stability by generating a restoring side force when the aircraft sideslips. The fin prevents unwanted yaw oscillations known as Dutch roll and serves as the mounting structure for the movable rudder. The size and position of the fin are designed based on the aircraft's directional stability requirements and are influenced by factors such as engine-out asymmetric thrust.}{catAE}

\dictentry{Flapping Hinge}{A hinge on a helicopter rotor hub that allows the blade to flap up and down, accommodating the dissymmetry of lift in forward flight by allowing the advancing blade to rise and the retreating blade to droop. The flapping hinge eliminates the strong rolling moment that would otherwise be imposed on the rotor hub by unequal lift distribution. On fully articulated rotors, the flapping hinge works in conjunction with the feathering and lead-lag hinges to provide full three-degree-of-freedom blade motion.}{catAE}

\dictentry{Flight Computer}{The central computing unit in a glass cockpit or avionics suite that integrates and processes data from navigation sensors, air data probes, and other avionics subsystems to drive flight displays and provide computations. The flight computer performs functions such as attitude computation, air data calculations, navigation solution generation, and flight plan management. Redundant flight computers operating in parallel ensure continuity of critical flight information in the event of a single unit failure.}{catAE}

\dictentry{Flight Director}{A cockpit instrument that computes and displays command guidance cues indicating what pitch and bank inputs the pilot should make to follow a selected flight path, such as a heading, altitude, or approach course. The flight director integrates data from the flight management system, navigation receivers, and air data computer to present steering bars on the attitude indicator. When coupled with the autopilot, the flight director becomes the visual representation of the autopilot's guidance commands.}{catAE}

\dictentry{Flight Management System}{An integrated avionics system that automates in-flight tasks by combining navigation, performance optimization, and autopilot control into a single platform. The FMS computes optimal flight profiles for fuel economy, manages the flight plan, and interfaces with the autopilot and autothrottle to fly the planned route with minimal pilot intervention. It accepts pilot-entered flight plans and performance data, then provides lateral and vertical guidance throughout all phases of flight from departure to arrival.}{catAE}

\dictentry{Flight Mechanics}{The branch of aerospace engineering that analyzes the forces acting on an aircraft in flight and the resulting motion, encompassing aerodynamics, propulsion, weight, and control. Flight mechanics is used to determine performance parameters such as range, endurance, climb rate, maneuverability, and stability characteristics across the flight envelope. Engineers use flight mechanics principles during the design phase to size wings, select engines, and establish the operating limitations published in the aircraft flight manual.}{catAE}

\dictentry{Flight Path}{The trajectory of the aircraft's center of gravity through three-dimensional space, defined by its ground track, altitude profile, and vertical speed along the route. The flight path is determined by the combination of aerodynamic forces, thrust, weight, and atmospheric conditions, and is influenced by pilot inputs and autopilot commands. Understanding the flight path is essential for collision avoidance, terrain clearance, and efficient fuel management throughout all phases of flight.}{catAE}

\dictentry{Flight Safety}{The application of engineering principles, operational procedures, and regulatory frameworks to prevent accidents and incidents during aircraft operation. Flight safety encompasses aircraft design and certification, crew training, maintenance programs, air traffic management, and operational procedures designed to minimize risk. The flight safety culture in aviation drives continuous improvement through accident investigation, incident reporting systems, and the implementation of corrective actions across the industry.}{catAE}

\dictentry{Flight Test}{The process of evaluating an aircraft's performance, handling qualities, and systems during actual flight conditions as part of certification or development. Flight test programs follow carefully planned test points covering the entire flight envelope including low speed, cruise, high altitude, and emergency conditions. Instrumented test aircraft record hundreds of parameters during each sortie, and results are compared against design predictions and certification requirements.}{catAE}

\dictentry{Flow Separation}{The detachment of the boundary layer from the surface of a wing or other aerodynamic body, which causes a dramatic loss of lift and a large increase in pressure drag. Flow separation occurs when the boundary layer encounters an adverse pressure gradient that it cannot overcome, causing the flow to reverse direction near the surface. At high angles of attack, progressive flow separation leads to aerodynamic stall, making it the primary phenomenon limiting the maximum lift capability of wings and control surfaces.}{catAE}

\dictentry{Flutter Test}{A flight test conducted to determine the speed at which aeroelastic flutter, a self-excited oscillation of wings, control surfaces, or other flexible structures, may occur. Flutter is potentially catastrophic because once initiated, oscillations can grow exponentially until structural failure results. Test pilots incrementally increase airspeed while monitoring vibration levels, and flutter margins are established to define safe operating speeds well below the onset point with an adequate safety margin.}{catAE}

\dictentry{Fly-by-Wire}{A flight control system in which pilot inputs are transmitted as electrical signals through a digital computer to the control surface actuators, replacing conventional mechanical linkages. Fly-by-wire systems offer weight savings, reduced maintenance, and the ability to implement control laws that enhance stability, provide envelope protection, and optimize maneuvering performance. The system includes redundant computers and sensors to ensure that no single failure can result in loss of control, and is standard on modern military and commercial aircraft.}{catAE}

\dictentry{Fowler Flap}{A high-lift trailing edge device that extends both aft and downward from the wing, simultaneously increasing the wing's effective chord length, camber, and surface area. The Fowler flap achieves the highest lift increment among common flap types, making it ideal for reducing approach and landing speeds on large transport aircraft. When retracted, the Fowler flap fits flush with the wing's lower surface, maintaining an efficient cruise airfoil shape.}{catAE}

\dictentry{Friction Stir Welding}{A solid-state joining process in which a rotating tool generates frictional heat to soften and plastically deform the material at the joint interface without melting it, producing a strong, defect-free weld. The process is particularly advantageous for joining aluminum alloys used in aerospace structures because it avoids the porosity and cracking issues associated with fusion welding. Friction stir welding is used on the SpaceX Falcon 9 rocket, the Boeing Delta IV, and in aircraft fuselage and fuel tank fabrication.}{catAE}

\dictentry{Fuel System}{The integrated assembly of tanks, pumps, valves, filters, fuel lines, and quantity indicators responsible for storing aviation fuel and delivering it to the engines at the correct pressure, temperature, and flow rate. Aircraft fuel systems must manage fuel transfer between multiple tanks to maintain center of gravity limits and ensure uninterrupted supply during all flight attitudes. Fuel system design must also address safety requirements including inerting of fuel tank vapors to prevent explosion in the event of a lightning strike or combat damage.}{catAE}

\dictentry{Fully Articulated Rotor}{A helicopter rotor system in which each blade is attached to the hub through three independent hinges allowing flapping, feathering, and lead-lag motion. The flapping hinge permits vertical blade movement to compensate for dissymmetry of lift, the feathering hinge enables pitch changes for collective and cyclic control, and the lead-lag hinge allows in-plane blade oscillation. This design offers the greatest mechanical complexity but provides the most flexible and responsive rotor dynamics.}{catAE}

\dictentry{Fuselage}{The central body of an aircraft that houses the flight crew, passengers, and cargo, and provides the structural framework to which the wings, tail, and landing gear are attached. Fuselage construction must balance aerodynamic streamlining with the internal volume requirements for pressurization, passenger comfort, and cargo loading. Modern fuselage designs use semi-monocoque construction with aluminum or composite skin supported by longitudinal stringers and transverse frames.}{catAE}

\dictentry{Generator}{A device that converts mechanical energy from a rotating shaft into electrical energy, typically driven by the aircraft engine through a gearbox or accessory drive. Aircraft generators produce alternating current at regulated voltage and frequency, with brushless designs preferred for their reliability and reduced maintenance. Modern systems use integrated drive generators that combine a constant-speed drive unit with the generator to maintain stable electrical output regardless of engine speed variations.}{catAE}


\dictentry{Glide Ratio}{The ratio of horizontal distance traveled to altitude lost during unpowered flight, numerically equal to the lift-to-drag ratio of the aircraft. A glider with a glide ratio of 40 to 1 can travel 40 kilometers forward while losing one kilometer of altitude in still air. The glide ratio varies with airspeed and is a key performance parameter for emergency descent planning and glider competition flying.}{catAE}

\dictentry{Glider}{A fixed-wing aircraft without an engine that relies entirely on rising air currents such as thermals, ridge lift, and wave lift to gain altitude and extend flight duration. Gliders are designed with extremely high lift-to-drag ratios, often exceeding 40 to 1, using long, narrow wings and streamlined composite fuselages to minimize drag. Modern gliders use variometers to detect rising air and optimize their flight path, with competitive flights covering hundreds of kilometers using only natural energy sources.}{catAE}

\dictentry{GNSS}{Global Navigation Satellite Systems, the collective term for satellite-based positioning constellations including the US GPS, Russian GLONASS, European Galileo, and Chinese BeiDou systems. GNSS receivers on aircraft calculate three-dimensional position and time by measuring signal travel times from multiple satellites, providing accuracy within meters for civil aviation applications. Multi-constellation receivers that track satellites from several systems simultaneously improve accuracy, availability, and integrity of the positioning solution.}{catAE}

\dictentry{GPS}{The Global Positioning System, a US-owned satellite constellation of at least 24 satellites orbiting at approximately 20,200 kilometers that provides continuous three-dimensional positioning, velocity, and time information to receivers worldwide. Aircraft GPS receivers determine position by calculating the time delay of signals from at least four satellites, achieving civil accuracy of about 3 to 5 meters. In aviation, GPS is used for en-route navigation, instrument approaches, and performance-based navigation procedures that increase airspace capacity and flexibility.}{catAE}

\dictentry{GPWS}{Ground Proximity Warning System, an avionics system that uses radar altimeter data and airframe parameters to warn the flight crew when the aircraft is in immediate danger of flying into terrain. GPWS provides distinctive aural warnings such as ``terrain, terrain, pull up'' and ``too low, terrain'' accompanied by visual alerts on the navigation display. Enhanced GPWS incorporates a forward-looking terrain awareness function using GPS terrain databases to provide warnings earlier than the basic radar-based system.}{catAE}

\dictentry{Gravity Turn}{A launch trajectory technique in which the vehicle initially ascends vertically and then gradually pitches over, allowing gravity to naturally curve the flight path toward horizontal for orbital insertion. The gravity turn minimizes aerodynamic loads on the vehicle during the high-dynamic-pressure phase of ascent and reduces the propellant required for the gravity loss component of the trajectory. This technique is used by virtually all launch vehicles from the Atlas and Falcon to the Space Shuttle and Ariane.}{catAE}

\dictentry{Great Circle Route}{The shortest path between two points on the surface of a sphere, which appears as a curved path on a flat map but represents the minimum geodesic distance on the curved Earth. Great circle navigation is fundamental to long-distance flight planning, as following a great circle route reduces total distance and fuel burn compared to flying along a constant compass heading. Airlines compute great circle routes between origin and destination pairs and then apply wind and airspace constraints to determine the optimal flight path.}{catAE}

\dictentry{Ground Effect Helicopter}{The enhanced aerodynamic efficiency that a helicopter rotor experiences when operating within approximately one rotor diameter of the ground, where the ground surface disrupts the formation of tip vortices and reduces induced drag. Ground effect allows a helicopter to hover at a lower power setting than it would require at higher altitudes, effectively increasing its hover ceiling near the surface. This phenomenon is exploited for hover-in-ground-effect operations near helipads, ship decks, and offshore platforms.}{catAE}

\dictentry{Ground Resonance}{A potentially destructive self-excited oscillation in helicopters where the rotor's out-of-plane motion couples with the airframe's natural frequencies, causing rapidly increasing vibrations that can destroy the aircraft within seconds. Ground resonance occurs when a fully articulated rotor is on the ground and one blade is displaced, causing an unbalanced rotor disc that excites the airframe's natural modes. The phenomenon is prevented through proper damping in the lead-lag hinge and by maintaining proper blade track and balance.}{catAE}

\dictentry{Ground Speed}{The horizontal speed of an aircraft measured relative to the Earth's surface, which differs from true airspeed due to the effect of wind. A headwind reduces ground speed below airspeed while a tailwind increases it, making ground speed the critical parameter for calculating time of arrival and fuel planning. Ground speed is displayed on GPS receivers and navigation instruments and is continuously updated as wind conditions change during flight.}{catAE}

\dictentry{Guidance}{The process of determining and maintaining the desired path or trajectory for a vehicle, whether an aircraft navigating between waypoints or a launch vehicle following a prescribed ascent profile. Guidance systems use sensors, algorithms, and control logic to compute the required steering commands that keep the vehicle on its intended path. In spacecraft, guidance is closely integrated with navigation and control in the GNC subsystem to execute precise orbital maneuvers and rendezvous operations.}{catAE}

\dictentry{Gust Alleviation}{An active control technique that uses accelerometers to detect sudden gust-induced load changes on the airframe and commands control surface deflections to counteract the resulting bending moments. By reducing peak gust loads, gust alleviation systems allow lighter structural designs and improve ride comfort for passengers. Modern systems on large transport aircraft use leading and trailing edge surfaces to actively counteract turbulence effects in real time.}{catAE}

\dictentry{Gyrocopter}{A rotorcraft, also called an autogyro, in which the main rotor is unpowered and spins freely in autorotation as air flows upward through the rotor disc during forward flight. A separate engine-driven propeller provides forward thrust, and the resulting airflow through the rotor generates lift. Gyrocopters cannot hover like helicopters but are simpler, lighter, and safer at low speeds because they are always in autorotation and are not subject to settling with power or vortex ring state.}{catAE}

\dictentry{H-V Curve}{The height-velocity diagram for helicopters, a chart that plots the combinations of altitude and airspeed from which a safe autorotation landing is possible following an engine failure. The curve defines a dead man's zone at low altitude and low airspeed where there is insufficient kinetic or potential energy to execute a successful autorotation. Pilots use the H-V diagram to avoid flight conditions that leave no margin for recovery in the event of power loss.}{catAE}

\dictentry{Handley Page Slat}{A leading edge aerodynamic device that automatically or manually extends forward from the wing's leading edge to delay stall at high angles of attack by creating a slot for high-pressure air to flow over the upper surface. The slat energizes the boundary layer, delaying flow separation and allowing the wing to achieve a higher maximum lift coefficient. Handley Page slats are found on many transport and general aviation aircraft and work in conjunction with trailing edge flaps for low-speed performance.}{catAE}

\dictentry{Hazard Analysis}{A systematic process of identifying potential hazards in an aerospace system, evaluating their likelihood and severity, and determining the mitigations needed to reduce risk to an acceptable level. Hazard analyses are performed throughout the product lifecycle from concept design through operations and include methods such as fault tree analysis, failure mode and effects analysis, and common cause analysis. The results feed into the safety assessment process required for airworthiness certification under standards such as ARP4761.}{catAE}

\dictentry{Head-Up Display}{A transparent display that projects critical flight information such as airspeed, altitude, heading, and flight director commands onto a combiner glass positioned directly in the pilot's forward field of view. The HUD allows the pilot to maintain outside visual reference while simultaneously monitoring flight instruments, significantly improving situational awareness during low-visibility operations. HUD systems are especially valuable during approaches, landings, and tactical military operations where heads-down time must be minimized.}{catAE}

\dictentry{Heat Shield}{A thermal protection system layer applied to the exterior of a spacecraft or re-entry vehicle that dissipates the extreme aerodynamic heating generated during atmospheric entry at orbital velocities. Heat shields use ablative materials that char and erode to carry heat away, or reusable ceramic tiles and tiles that radiate thermal energy back to the atmosphere. The design must withstand temperatures exceeding 1,600 degrees Celsius while protecting the underlying structure and crew compartment from thermal damage.}{catAE}

\dictentry{Height-Velocity Diagram}{A graphical representation used in helicopter operations that shows the combinations of altitude and airspeed from which a successful autorotation landing can be performed following an engine failure. The diagram defines the avoid area at low altitude and low airspeed where recovery is not possible due to insufficient rotor energy. Pilots consult the height-velocity diagram before takeoff and during flight to ensure they remain in regions where a safe power-off landing can be accomplished.}{catAE}

\dictentry{Helicopter}{A rotary-wing aircraft that achieves lift, thrust, and directional control through one or more horizontally rotating rotor systems powered by one or more turbine or piston engines. Unlike fixed-wing aircraft, helicopters can hover, take off and land vertically, and fly in any direction, making them indispensable for search and rescue, offshore transport, military operations, and emergency medical services. The complex aerodynamics of helicopter flight involve phenomena such as dissymmetry of lift, vortex ring state, and ground resonance that require specialized pilot training.}{catAE}

\dictentry{High-Lift Device}{Any aerodynamic mechanism on a wing that increases its maximum lift coefficient for low-speed flight during takeoff and landing, including flaps, slats, leading edge devices, and spoilers used as lift dumpers. High-lift devices allow aircraft to operate from shorter runways by reducing the minimum speed at which the wing can generate sufficient lift to support the aircraft's weight. They are retracted during cruise to maintain the clean wing's efficient aerodynamic profile.}{catAE}

\dictentry{Hohmann Transfer}{An efficient two-impulse orbital transfer between two circular orbits in the same plane, consisting of a burn to enter an elliptical transfer orbit followed by a second circularization burn at the target altitude. The Hohmann transfer uses the minimum propellant of any two-impulse transfer between coplanar circular orbits and was first described by German engineer Walter Hohmann in 1925. It is widely used for interplanetary transfers and orbit-raising maneuvers, with typical transfer times depending on the size of the orbit change.}{catAE}



\dictentry{Horizontal Stabilizer}{The fixed horizontal tail surface mounted at the rear of the aircraft that provides longitudinal stability by generating a downward or upward aerodynamic force to balance the pitching moment created by the wing and fuselage. The horizontal stabilizer also serves as the mounting structure for the elevator or stabilator control surface. Its size and position are determined by the required static margin, which ensures the aircraft naturally tends to return to its trimmed angle of attack after a disturbance.}{catAE}

\dictentry{Horn Balance}{A forward-projecting extension of a control surface ahead of the hinge line that increases aerodynamic balance by exposing a portion of the surface to the airstream in front of the pivot point. The horn balance reduces the control force required by the pilot by generating a moment that partially counteracts the hinge moment produced by the main portion of the surface behind the hinge. Careful design is required to avoid overbalance conditions that could lead to control surface flutter or instability.}{catAE}

\dictentry{Hot Staging}{A rocket staging technique in which the upper stage engine ignites while still attached to the lower stage, before separation occurs, ensuring continuous thrust and preventing propellant settling in microgravity. Hot staging is more mechanically complex than cold staging because it requires venting the interstage to prevent pressure buildup, but it eliminates the coast phase and reduces gravity losses. SpaceX's Starship uses hot staging with a vented interstage ring to maintain thrust continuity during stage separation.}{catAE}

\dictentry{Hover}{A flight condition in which a helicopter or VTOL aircraft maintains a stationary position in the air at zero ground speed, with lift exactly equal to weight and thrust directed vertically upward. Hovering is one of the most power-demanding flight conditions for a helicopter because the rotor must generate sufficient lift without the benefit of translational lift. The hover performance is characterized by the hover ceiling, which decreases with increasing altitude, temperature, and aircraft weight.}{catAE}

\dictentry{Hybrid-Electric Propulsion}{A propulsion architecture that combines conventional gas turbine or piston engines with electric motors and batteries to achieve improved fuel efficiency, reduced emissions, or enhanced operational flexibility. In a parallel hybrid, both the thermal engine and electric motor can drive the propulsor, while in a series hybrid the gas turbine acts as a generator to power electric motors. Hybrid-electric designs are being explored for regional aircraft and urban air mobility vehicles as an intermediate step toward fully electric flight.}{catAE}

\dictentry{Hydraulic System}{A power transmission system that uses pressurized fluid, typically mineral-based or synthetic hydraulic fluid, to actuate flight controls, landing gear, brakes, and other high-force mechanisms throughout the aircraft. Hydraulic systems provide high power-to-weight ratios and precise control authority, with operating pressures typically at 3,000 psi in commercial aircraft and up to 5,000 psi in military systems. Most transport aircraft have two or three independent hydraulic circuits with cross-feed capability to ensure control system availability after a single system failure.}{catAE}

\dictentry{Hydrogen Aircraft}{An aircraft that uses hydrogen as its primary fuel, either burned directly in modified gas turbines or converted to electricity through fuel cells, offering the potential for zero-carbon flight since hydrogen combustion produces only water vapor. Liquid hydrogen must be stored at cryogenic temperatures of minus 253 degrees Celsius in heavily insulated tanks, which presents significant design challenges for aircraft volume and weight. Several manufacturers and startups are developing hydrogen-powered concepts for short- to medium-range aviation as part of the industry's decarbonization strategy.}{catAE}

\dictentry{Hypersonic Vehicle}{A craft capable of sustained flight at speeds exceeding Mach 5, where extreme aerodynamic heating, plasma formation, and compressibility effects create unique engineering challenges. Hypersonic vehicles require specialized thermal protection systems, scramjet or dual-mode propulsion engines, and heat-resistant materials such as ceramic matrix composites and ultra-high-temperature ceramics. Applications include high-speed transport, space access vehicles, and military reconnaissance platforms.}{catAE}

\dictentry{ICAO}{The International Civil Aviation Organization, a specialized United Nations agency established in 1944 that sets international standards and recommended practices for civil aviation safety, security, efficiency, and environmental protection. ICAO coordinates with 193 member states to harmonize regulations covering aircraft certification, airworthiness, pilot licensing, air traffic management, and airport operations. The organization issues Standards and Recommended Practices that form the basis of each member state's national aviation regulations.}{catAE}

\dictentry{Ice Protection System}{A system that prevents or removes ice accumulation on critical aircraft surfaces such as wings, engine inlets, pitot tubes, and tail surfaces where ice can degrade aerodynamic performance and engine operation. Common ice protection methods include thermal anti-ice using engine bleed air, electrothermal de-ice boots that cycle to break ice bonds, and fluid-based systems that apply anti-icing chemicals to leading edges. Ice protection systems are essential for all-weather operations and are rigorously tested and certified for the aircraft's icing envelope.}{catAE}

\dictentry{Indicated Airspeed}{The speed read directly from the aircraft's airspeed indicator, derived from the difference between pitot (ram) pressure and static pressure without any corrections. Indicated airspeed is affected by instrument errors, position errors from the placement of pitot and static probes, and compressibility effects at high speeds. Pilots use indicated airspeed for operational reference because it directly relates to aerodynamic forces and is the basis for V-speeds and performance calculations.}{catAE}

\dictentry{INS}{An Inertial Navigation System, a self-contained navigation aid that uses accelerometers and gyroscopes to measure the aircraft's acceleration and rotation, integrating these measurements over time to compute position, velocity, and attitude without external references. INS is immune to signal jamming and does not require satellite signals, making it valuable for military and long-range operations. However, its position accuracy degrades over time due to sensor drift, so it is typically periodically updated from GPS or other external navigation sources.}{catAE}

\dictentry{Internal Balance}{A sealed chamber built into a control surface forward of the hinge line that uses trapped air or a vacuum to reduce the aerodynamic hinge moment by equalizing pressure on both sides of the surface. Unlike external balance tabs, internal balance systems are fully enclosed and do not create additional drag or require external linkages. They are commonly used on combat aircraft control surfaces where low drag and clean aerodynamic profiles are essential.}{catAE}

\dictentry{Inverter}{A power electronics device that converts direct current from the aircraft's DC electrical system or batteries into alternating current at the required voltage and frequency to power AC loads. Inverters are essential for operating equipment designed for AC power, including avionics, lighting, and galley equipment, when only DC power is available from batteries or emergency sources. Modern inverters use solid-state switching devices for high efficiency and are built with redundancy to ensure uninterrupted AC power supply.}{catAE}

\dictentry{Jet Flap}{A high-lift device in which engine exhaust or a dedicated blower ejects high-velocity air over a hinged flap at the trailing edge, dramatically increasing circulation and lift coefficient. The jet flap effectively turns the entire wing into a circulation control surface, generating lift coefficients far exceeding those of conventional flaps. Though highly effective, the complexity and weight penalties of the required ducting have limited its use primarily to research and STOL military aircraft.}{catAE}

\dictentry{Kalman Filter}{A recursive mathematical algorithm that estimates the true state of a dynamic system by fusing predictions from a mathematical model with noisy measurements from multiple sensors, optimally weighting each input based on its estimated uncertainty. The Kalman filter is fundamental to modern navigation, guidance, and control systems in both aircraft and spacecraft. Extended and unscented variants handle the nonlinearities common in real aerospace applications such as orbit determination and attitude estimation.}{catAE}

\dictentry{Kevlar}{An aramid synthetic fiber with exceptional tensile strength and impact resistance, five times stronger than steel on a weight-for-weight basis, used in aerospace for ballistic protection, structural reinforcement, and containment applications. Kevlar is used in helicopter floor armor, turbine engine fan blade containment rings, and pressure vessel liners to prevent catastrophic failure from fragment impact. Its low density and high energy absorption characteristics make it ideal for protective structures where weight is critical.}{catAE}

\dictentry{Krueger Flap}{A leading edge high-lift device that hinges from the lower surface of the wing and extends forward and downward to increase camber and delay stall at high angles of attack. Unlike slats which create a slot, Krueger flaps modify the leading edge profile directly and are often used on the inboard wing sections of transport aircraft where Fowler flaps are not practical. Boeing 737 aircraft use Krueger flaps on the wing leading edge to complement the slotted flap system.}{catAE}

\dictentry{Landing Gear}{The undercarriage assembly of an aircraft that supports the vehicle during ground operations including taxiing, takeoff roll, and landing, consisting of wheels, tires, shock absorbers, brakes, and retraction mechanisms. Landing gear design must absorb the kinetic energy of descent rates of several meters per second while supporting the full aircraft weight and transmitting loads into the primary structure. Most retractable landing gear systems are actuated hydraulically or electrically and include position indicators and warning systems to ensure the gear is down and locked before landing.}{catAE}



\dictentry{Lead-Lag Hinge}{A hinge on a helicopter rotor hub that allows the blade to oscillate in the plane of rotation, accommodating the Coriolis forces generated as the blade flaps up and down during each revolution. The lead-lag motion is essential for preventing excessive hub loads and structural fatigue in fully articulated rotors. Dampers are typically installed at the lead-lag hinge to suppress ground resonance and prevent destructive oscillations during ground operations.}{catAE}

\dictentry{Leading Edge Device}{An aerodynamic device mounted on the leading edge of a wing that improves low-speed lift characteristics by modifying the airfoil shape or energizing the boundary layer. Types include fixed slots, automatic and handley page slats, Krueger flaps, and leading edge cuffs, each offering different trade-offs between high-lift gain, drag penalty, and mechanical complexity. Leading edge devices are critical for maintaining adequate stall margins and reducing approach speeds during landing.}{catAE}

\dictentry{Linkage}{The system of mechanical connections including pushrods, cables, pulleys, bellcranks, and control rods that transmits the pilot's control inputs from the cockpit to the flight control surfaces. Linkage systems are designed with minimal free play and backlash to provide precise and predictable control response. In fly-by-wire aircraft, traditional mechanical linkages are replaced by electrical signal pathways, though redundant mechanical connections may be retained as backups in some designs.}{catAE}

\dictentry{Load Alleviation}{An active control strategy that uses sensors to detect structural loads in real time and commands control surface deflections to redistribute and reduce peak bending and torsional moments on the airframe. By reducing the maximum loads experienced during gusts and maneuvers, load alleviation allows designers to use lighter structural members without sacrificing fatigue life. Modern load alleviation systems are integrated into the fly-by-wire control laws of large transport and military aircraft.}{catAE}

\dictentry{Longeron}{A major longitudinal structural member of the fuselage or tail boom that runs along the length of the airframe and carries axial loads including bending, tension, and compression. Longeronss are typically heavy-gauge extruded or machined aluminum or composite sections located at the corners or critical load paths of the fuselage cross-section. They work in conjunction with transverse frames and stringers to form the semi-monocoque structural framework that carries all flight loads.}{catAE}



\dictentry{Main Rotor}{The primary lifting rotor assembly of a helicopter, consisting of two or more blades attached to a rotor hub that is driven by the engine through a transmission. The main rotor generates both the lift required to support the aircraft's weight and the thrust for forward flight when tilted by the cyclic control. Main rotor systems are classified by their blade attachment method as fully articulated, semi-rigid, or rigid, each offering different performance and handling characteristics.}{catAE}

\dictentry{Mass Balance}{Weights attached to control surfaces ahead of their hinge line to move the center of gravity of the surface forward to or ahead of the hinge axis, preventing aerodynamic flutter. Mass balancing ensures that inertial forces during oscillation do not add energy to a potential flutter mode, which could lead to catastrophic structural failure. The mass balance condition is verified during certification and must be maintained throughout the aircraft's operational life.}{catAE}

\dictentry{Master Warning}{A prominent visual alert indicator in the cockpit, typically an illuminated red or amber light, that draws the pilot's attention to an abnormal condition requiring immediate action or awareness. The master warning light is accompanied by specific system alerts and caution messages that identify the exact nature of the malfunction. When activated, the pilot acknowledges the warning by pressing the master warning button, which silences the aural alert while the condition persists.}{catAE}


\dictentry{MIL-STD-1553}{A United States military standard defining the requirements for a serial data bus communication system used in military aircraft, spacecraft, and ground vehicles for deterministic, reliable data exchange between avionics components. The standard specifies a dual-redundant twisted-pair bus architecture, a command-response protocol with fixed message formats, and error detection through Manchester encoding and parity checking. MIL-STD-1553 has been the backbone of military avionics data communication for decades and has been adopted by numerous NATO allies.}{catAE}

\dictentry{Mixture Control}{A cockpit control that adjusts the ratio of fuel to air being delivered to a piston engine, allowing the pilot to optimize combustion efficiency for different altitudes and power settings. At higher altitudes where air density decreases, the mixture must be leaned to avoid an overly rich condition that wastes fuel and fouls spark plugs. Modern fuel-injected engines with automatic mixture control eliminate the need for manual mixture adjustment, but carbureted engines still require pilot-managed mixture control.}{catAE}

\dictentry{Mode Control Panel}{The cockpit panel, typically located on the glare shield, through which the pilot selects autopilot modes such as heading hold, altitude select, vertical speed, and approach coupling. The MCP provides the primary interface between the pilot and the autopilot flight director system, displaying the selected modes and current values for each parameter. Mode annunciations on the primary flight display confirm which modes are active and armed, helping the pilot maintain awareness of the autopilot's behavior.}{catAE}

\dictentry{Monocoque}{A structural design in which the outer skin carries all or most of the primary loads, eliminating the need for an internal framework of longerons and ribs, similar to the shell of an egg. Monocoque construction provides excellent strength-to-weight ratio and smooth aerodynamic surfaces but is vulnerable to localized damage that can propagate as buckling across the unsupported skin panels. Pure monocoque is rare in modern aircraft and has largely been replaced by semi-monocoque designs that add internal reinforcement.}{catAE}

\dictentry{Morphing Wing}{A wing structure capable of changing its shape in flight to optimize aerodynamic performance across different flight conditions, including span, camber, twist, and sweep adjustments. Morphing wing technology uses flexible skins, compliant mechanisms, and shape memory alloys to achieve smooth, continuous deformations without the gaps and steps of conventional hinged control surfaces. Research in morphing wings promises significant improvements in fuel efficiency by allowing the wing to always operate at its aerodynamic optimum.}{catAE}

\dictentry{MTBF}{Mean Time Between Failures, a reliability metric representing the average operating time expected between consecutive failures of a repairable component or system. MTBF is calculated by dividing total operating hours by the number of failures during a specified period and is used to schedule preventive maintenance, predict spare parts requirements, and assess system availability. In aerospace, MTBF values are critical for certifying the reliability of essential systems and establishing maintenance program intervals.}{catAE}

\dictentry{Multi-Function Display}{A configurable electronic display screen in a glass cockpit that can present multiple types of information including navigation maps, weather radar, engine parameters, and system synoptic pages at the pilot's selection. MFDs replaced multiple dedicated instruments with a single versatile display, reducing cockpit complexity and weight. The pilot selects the desired page or format using display select keys, allowing the same screen to serve different purposes during different phases of flight.}{catAE}

\dictentry{Navigation}{The science and practice of determining an aircraft's position, planning a course to a destination, and monitoring progress along the planned route throughout a flight. Modern aviation navigation integrates multiple systems including GPS, inertial reference units, VOR and DME ground-based aids, and the flight management system to provide continuous position updates. Pilotage, dead reckoning, and electronic navigation methods are all taught as complementary skills for safe and efficient flight operations.}{catAE}

\dictentry{Navigation Display}{A cockpit display that presents navigation information including current position, planned route, nearby waypoints, navaids, airspace boundaries, and weather radar returns in a moving map format. The navigation display is driven by the flight management system and GPS, providing the pilot with situational awareness of the aircraft's position relative to the flight plan and surrounding traffic. During instrument approaches, the display shows the approach path, glideslope, and localizer deviation for precision guidance.}{catAE}

\dictentry{Nose Gear}{The forward landing gear assembly of a tricycle-gear aircraft, consisting of one or two wheels mounted on a steerable strut that supports the nose section of the fuselage during ground operations. The nose gear provides directional steering during taxi and absorbs landing loads in nose-first contact scenarios. It typically retracts forward or aft into the fuselage and includes steering actuators linked to the pilot's rudder pedals or a tiller wheel.}{catAE}

\dictentry{NOTAR}{No Tail Rotor, a helicopter anti-torque system developed by McDonnell Douglas that eliminates the conventional tail rotor by using a pressurized air system and variable-incidence tail boom to provide directional control. NOTAR reduces noise, eliminates the danger of tail rotor strikes on personnel and objects, and simplifies maintenance. The system works by circulating compressed air through slots in the tail boom to exploit the Coanda effect, directing airflow to generate the necessary anti-torque and control forces.}{catAE}

\dictentry{Orbit}{The curved gravitational path of an object around a celestial body, maintained by the balance between the object's inertia and the gravitational attraction of the central body. Orbits can be circular, elliptical, parabolic, or hyperbolic depending on the object's velocity relative to escape speed. The shape and size of an orbit are described by six classical orbital elements including semi-major axis, eccentricity, inclination, and right ascension of the ascending node.}{catAE}

\dictentry{Orbital Inclination}{The angle between the plane of an orbit and a reference plane, typically the Earth's equatorial plane, measured at the ascending node where the satellite crosses the reference plane heading northward. An inclination of zero degrees indicates an equatorial orbit, while 90 degrees defines a polar orbit, and values between provide varying degrees of latitudinal coverage. The inclination determines the maximum latitude over which a satellite passes directly overhead and is a fundamental parameter in orbit design and launch vehicle selection.}{catAE}


\dictentry{Parking Brake}{A mechanical locking device that holds the aircraft brakes in the applied position to prevent the aircraft from rolling while parked on the ground. The parking brake is typically engaged by applying wheel brakes and then setting a mechanical lock that maintains hydraulic pressure in the brake cylinders. Pilots set the parking brake before engine shutdown and release it before taxi, and it is also used as a backup during ground maintenance operations.}{catAE}

\dictentry{Payload}{The useful cargo delivered by a launch vehicle to its target orbit, which may include satellites, space station resupply modules, scientific instruments, or crewed spacecraft. Payload capacity to a specified orbit is the primary metric for comparing launch vehicles and directly determines mission capability and cost per kilogram to orbit. The payload is enclosed within the payload fairing during ascent and is deployed or docked once the vehicle reaches the target orbital parameters.}{catAE}

\dictentry{PID Controller}{A feedback control algorithm that calculates an output signal based on the sum of three terms proportional to the error, the integral of the error, and the derivative of the error. The proportional term reduces current error, the integral term eliminates steady-state offset, and the derivative term damps oscillations and improves stability. PID controllers are ubiquitous in aerospace applications including autopilot hold modes, engine fuel control, and attitude regulation due to their simplicity and effectiveness.}{catAE}

\dictentry{Pneumatic System}{A system that uses compressed air or other gases to power aircraft subsystems including brake actuation, door operation, wing de-icing boots, and water waste valve control. Pneumatic systems are powered by engine bleed air or dedicated compressors and operate at pressures typically between 30 and 150 psi. They offer advantages in weight, simplicity, and tolerance to contamination compared to hydraulic systems, though they deliver less force for a given actuator size.}{catAE}


\dictentry{Power Controls}{Fully powered flight control systems in which pilot inputs are transmitted directly to hydraulic or electric actuators that move the control surfaces, with no mechanical connection between the cockpit and the surfaces. In these systems, the pilot's controls serve only as input transducers, and all force to move the surfaces is provided by the actuators. Power controls offer significant weight savings and allow envelope protection and load limiting but require robust redundancy to ensure control authority after a single actuator or power source failure.}{catAE}

\dictentry{Power Required Curve}{A graphical plot showing the total power needed to maintain level flight at various airspeeds, shaped like a U with minimum power at the best endurance speed and minimum drag at the best glide speed. The curve reflects the combined effects of parasite drag, which increases with speed, and induced drag, which decreases with speed. Pilots and designers use this curve to identify optimal cruise speeds, range speeds, and climb performance across the flight envelope.}{catAE}

\dictentry{Pressurization System}{A system that maintains a controlled atmospheric pressure inside the aircraft cabin to ensure passenger and crew comfort and safety during flight at altitudes where the outside air is too thin for normal breathing. The system uses engine bleed air or electric compressors to supply air, while outflow valves regulate the rate of air release to control cabin altitude. The structural design of the fuselage must withstand the differential pressure, which creates significant hoop and longitudinal stresses in the skin and frames.}{catAE}

\dictentry{Primary Control}{The main aerodynamic surfaces responsible for controlling an aircraft's rotation about its three axes: elevators for pitch, ailerons for roll, and the rudder for yaw. Primary control surfaces are the most critical flight control elements, and their failure or malfunction can severely compromise the pilot's ability to safely maneuver the aircraft. Certification requirements mandate that primary controls meet the highest reliability and redundancy standards, with backup systems available for transport category aircraft.}{catAE}

\dictentry{Primary Flight Display}{The principal electronic display in a glass cockpit that presents the aircraft's attitude, airspeed, altitude, vertical speed, heading, and flight director cues in an integrated format. The PFD replaces the traditional six-pack of analog instruments with a single high-resolution screen that provides enhanced situational awareness through color coding, trend vectors, and alerting. It is driven by data from the air data computer, inertial reference system, and GPS receivers.}{catAE}

\dictentry{Propeller Control}{The system that manages propeller RPM and blade pitch angle to optimize engine performance and fuel efficiency across different flight conditions, consisting of a governor, pitch change mechanism, and oil transfer system. In constant-speed propellers, the governor automatically adjusts blade pitch to maintain a selected RPM regardless of airspeed or power setting. Feathering capability allows the blades to be turned edge-on to the airflow in an engine failure scenario, reducing drag on multi-engine aircraft.}{catAE}

\dictentry{Pulsejet}{A simple jet engine that operates through intermittent combustion cycles, in which air and fuel enter a combustion chamber through one-way valves, ignite to produce a high-pressure exhaust pulse, and then expel spent gases before the next cycle begins. Pulsejets are mechanically simple and inexpensive but generate high noise levels and vibration, limiting their use to expendable applications such as target drones and cruise missiles. The V-1 flying bomb of World War II was the most well-known operational pulsejet-powered vehicle.}{catAE}

\dictentry{Pushrod}{A rigid rod used in mechanical flight control systems to transmit control input forces from the cockpit controls or bellcranks to the control surface horns or actuators. Pushrods are typically made from steel tubing or aluminum alloy and operate in compression and tension along their axis. They are preferred over cables in applications requiring precise positional control and are common in light aircraft and the inboard sections of larger aircraft control systems.}{catAE}

\dictentry{Quadrant}{A sector-shaped control lever in the cockpit used for adjusting engine power or propeller settings, where the arc of travel corresponds to the range of available settings. The throttle quadrant typically houses multiple levers for power control, propeller RPM, and mixture, each marked with detents for standard positions such as idle, cruise, and full power. Quadrant controls are designed for positive tactile identification so pilots can adjust settings without looking away from instruments or outside references.}{catAE}

\dictentry{Qualification Test}{A formal test procedure conducted to demonstrate that a specific design meets all of its predefined performance, environmental, and durability requirements before certification approval. Qualification tests typically include structural proof loading, environmental stress screening for temperature and vibration, electromagnetic compatibility testing, and accelerated life cycle testing. The results are documented in a qualification test report that forms part of the certification data package submitted to the aviation authority.}{catAE}

\dictentry{Radar Altimeter}{An instrument that measures the aircraft's height above the terrain directly below by transmitting a radio signal downward and timing the reflected return, providing accurate altitude data from the surface up to approximately 2,500 feet. Unlike barometric altimeters, the radar altimeter provides true height above ground level regardless of atmospheric pressure or temperature conditions. It is essential for low-level flight, approach and landing guidance, terrain awareness systems, and automatic landing systems.}{catAE}





\dictentry{Rectifier}{A device that converts alternating current to direct current, essential in aircraft electrical systems where both AC power from generators and DC power for avionics, batteries, and DC-operated equipment are required. Rectifiers use semiconductor diodes arranged in bridge or center-tap configurations to produce smooth DC output from the AC input. Modern aircraft power systems incorporate solid-state rectifiers with thermal management and redundancy to ensure reliable DC power conversion.}{catAE}

\dictentry{Redundancy}{The incorporation of duplicate or backup components and systems in an aircraft design to ensure continued safe operation in the event of a single failure or malfunction. Redundancy levels are specified by certification regulations based on the severity of the failure condition, with catastrophic failures requiring at least three independent redundant paths. Common examples include dual hydraulic systems, triple-redundant flight computers, and multiple generators on separate engine drives.}{catAE}

\dictentry{Reliability}{The probability that an aerospace component or system will perform its intended function without failure for a specified period under stated operating conditions. Reliability is quantified through statistical analysis of test data and field experience, and is expressed as a probability or as mean time between failures. Design techniques for achieving required reliability levels include redundancy, derating, parts screening, and fault tolerance, all of which are central to the certification of flight-critical systems.}{catAE}

\dictentry{Resin Transfer Molding}{A closed-mold composite manufacturing process in which dry fiber preforms are placed in a mold and liquid resin is injected under pressure to impregnate the fibers before curing. RTM produces parts with excellent surface finish on both sides, controlled fiber volume fraction, and repeatable dimensional accuracy. The process is suitable for medium-to-high production volumes of complex-shaped aerospace components such as structural brackets, engine cowlings, and fairings.}{catAE}

\dictentry{Retreating Blade Stall}{An aerodynamic condition on a helicopter's retreating rotor blade where the blade's angle of attack exceeds the stall limit due to the combination of flapping-induced reduced airspeed and pilot cyclic input. This phenomenon limits the maximum forward speed of a conventional helicopter, as exceeding the retreating blade stall speed results in vibration, loss of lift, and potential loss of control. Dissymmetry of lift compensation through blade flapping reduces but cannot fully eliminate this fundamental speed limitation.}{catAE}

\dictentry{Reusable Launch Vehicle}{A rocket or spacecraft system designed to be recovered and refurbished for multiple flights, dramatically reducing the cost of access to space compared to expendable launch vehicles. Reusability strategies range from propulsive landing of first-stage boosters as demonstrated by the SpaceX Falcon 9, to full vehicle reusability with winged return like the retired Space Shuttle. The economic viability of reusable launch vehicles depends on rapid turnaround times, minimal refurbishment requirements, and high flight rates to amortize development costs.}{catAE}

\dictentry{Rib}{A chordwise structural member of a wing or tail surface that defines the airfoil cross-sectional shape and transfers aerodynamic loads from the skin to the spars. Ribs are typically made from aluminum sheet metal or composite layups and are spaced at regular intervals along the span to maintain the aerodynamic contour. Ribs also serve as attachment points for control surface hinges, fuel tank baffles, and leading and trailing edge devices.}{catAE}

\dictentry{Rigid Rotor}{A helicopter rotor system that has no mechanical hinges at the hub, relying instead on the inherent flexibility of the composite blade root and hub structure to accommodate flapping, feathering, and lead-lag motion. Rigid rotors offer simpler hub construction, reduced maintenance, and more responsive handling characteristics compared to articulated systems. The Eurocopter Tiger and the MBB Bo 105 are notable examples of helicopters that use rigid or bearingless rotor systems.}{catAE}

\dictentry{Risk Assessment}{A systematic process of identifying potential hazards, evaluating the likelihood and consequences of their occurrence, and determining appropriate mitigation measures to reduce risk to an acceptable level. In aerospace, risk assessments are performed at every stage from concept development through operations and retirement, using tools such as fault tree analysis, failure modes and effects analysis, and probabilistic risk assessment. The results inform design decisions, operational procedures, maintenance programs, and regulatory requirements.}{catAE}

\dictentry{Rotor Brake}{A mechanical braking device installed on the main rotor mast or transmission output shaft to decelerate and stop the rotor after engine shutdown and to prevent windmilling during ground operations. Rotor brakes use carbon or steel friction surfaces and are essential for reducing the time between landing and safe crew/passenger egress. On some helicopters the rotor brake also serves as a rotor lock during shipping or storage to prevent blade movement.}{catAE}

\dictentry{Rotor Head}{The mechanical assembly at the top of the rotor mast that connects the main rotor blades to the drive system and provides the necessary degrees of freedom for blade pitch, flapping, and lead-lag motion. The rotor head design determines whether the helicopter uses a fully articulated, semi-rigid, or rigid rotor system, and must transmit all aerodynamic loads from the blades into the airframe. Rotor head components are among the most highly stressed and fatigue-critical parts in any helicopter.}{catAE}

\dictentry{Rotorcraft}{A category of heavier-than-air aircraft that generates lift and thrust through one or more rotating wings or rotors, including helicopters, gyrocopters, and compound rotorcraft designs. Rotorcraft are distinguished from fixed-wing aircraft by their ability to hover, take off and land vertically, and maneuver in any direction at low speeds. The complex unsteady aerodynamics of rotating blades create unique flight challenges including dissymmetry of lift, vortex interactions, and ground effect that require specialized design and piloting techniques.}{catAE}



\dictentry{Secondary Control}{Aerodynamic surfaces and devices that supplement the primary controls to improve flight characteristics, reduce pilot workload, or enhance performance, including trim tabs, flaps, spoilers, and speed brakes. Secondary controls are not essential for maintaining basic flight control but significantly improve handling qualities, fuel efficiency, and operational capability. Their failure is classified at a lower severity level than primary control failures in the certification safety assessment process.}{catAE}

\dictentry{Semi-Monocoque}{A stressed-skin structural design in which loads are shared between the outer skin, longitudinal stringers, transverse frames, and bulkheads, providing multiple load paths for fail-safe behavior. Semi-monocoque construction is the dominant structural approach for modern aircraft fuselages and wings, offering an excellent balance of strength, stiffness, weight, and damage tolerance. Unlike pure monocoque structures, the internal reinforcement allows semi-monocoque designs to sustain localized damage without catastrophic structural failure.}{catAE}

\dictentry{Semi-Rigid Rotor}{A helicopter rotor system that uses a single teetering hinge to allow the rotor disc to tilt as a whole, combining the flapping and lead-lag functions into one degree of freedom at the hub. This simpler design reduces maintenance and component count compared to fully articulated rotors but limits the independent blade flapping capability. Semi-rigid rotors are commonly found on two-bladed helicopters such as the Bell JetRanger and Robinson R22 series.}{catAE}

\dictentry{Sensor Fusion}{The process of combining data from multiple sensors with different characteristics, such as GPS, inertial measurement units, air data probes, and vision systems, to produce a state estimate that is more accurate and reliable than any single sensor could provide. Sensor fusion algorithms such as Kalman filters weight each input based on its known accuracy and timing to produce optimal combined estimates. In aerospace, sensor fusion is essential for navigation, attitude determination, and autonomous flight where no single sensor provides complete or sufficiently accurate information.}{catAE}

\dictentry{Settling With Power}{A dangerous helicopter flight condition, also known as vortex ring state, in which the aircraft descends into its own rotor downwash, causing a recirculation of air that dramatically reduces rotor efficiency and lift. This condition typically occurs during steep descents at low airspeeds with power applied, and can lead to an uncontrolled descent rate even at full power. Recovery requires the pilot to pitch the nose forward and accelerate out of the descending vortex ring, trading vertical speed for forward airspeed.}{catAE}

\dictentry{Shape Memory Alloy}{An alloy, typically nickel-titanium, that can be deformed at low temperature and return to its original shape when heated above a specific transformation temperature, producing useful mechanical work in the process. Shape memory alloys enable compact, lightweight actuators and morphing structures without motors, gears, or hydraulic actuators. In aerospace, they are used for variable geometry engine inlets, deployable space structures, vibration damping systems, and self-healing seals.}{catAE}

\dictentry{Sideslip Angle}{The angle between the aircraft's longitudinal axis and the relative wind direction in the horizontal plane, indicating whether the aircraft is flying nose-left or nose-right of its actual flight path. A nonzero sideslip angle generates lateral aerodynamic forces that activate the vertical stabilizer's directional stability, tending to realign the nose with the relative wind. Pilots use controlled sideslip inputs during crosswind landings and single-engine operations, and the sideslip indicator is essential for coordinated flight.}{catAE}

\dictentry{Skin}{The outer covering of an aircraft that transfers aerodynamic loads to the internal structure and provides the aerodynamic surface shape, typically made from aluminum alloy or composite material panels fastened to ribs, stringers, and frames. The skin must resist buckling under compressive loads, withstand fatigue from pressurization cycles, and maintain its surface quality for efficient airflow. In semi-monocoque designs, the skin carries a significant portion of the flight loads in addition to its aerodynamic function.}{catAE}

\dictentry{Slotted Flap}{A trailing edge flap that creates a narrow gap or slot between the flap and the wing trailing edge when extended, allowing high-pressure air from below the wing to flow over the upper surface of the flap and energize the boundary layer. This slot effect delays flow separation on the deflected flap, producing higher lift coefficients than a plain flap of equivalent deflection angle. Slotted flaps are the most common type used on transport and business aircraft due to their favorable lift-to-drag ratio at approach speeds.}{catAE}





\dictentry{Sounding Rocket}{A suborbital research rocket that carries scientific instruments into the upper atmosphere and near-space environment for brief measurement durations of several minutes before falling back to Earth. Sounding rockets follow a ballistic trajectory that provides microgravity or reduced gravity conditions and access to altitudes between 50 and 1,500 kilometers, filling the gap between balloon-borne experiments and orbital missions. They are widely used by universities and research institutions for atmospheric science, astronomy, and technology demonstration at relatively low cost.}{catAE}

\dictentry{Space Debris}{Non-functional artificial objects in Earth orbit including defunct satellites, spent rocket stages, fragmentation fragments, and mission-related hardware that pose collision hazards to operational spacecraft. The Kessler Syndrome describes a scenario in which cascading collisions between debris objects generate increasing numbers of fragments, potentially rendering certain orbital regimes unusable. Space agencies track debris objects larger than 10 centimeters and implement collision avoidance maneuvers for operational satellites and crewed spacecraft.}{catAE}

\dictentry{Space Plasma}{Hot ionized gas composed of free electrons and ions that fills much of the space environment, including the solar wind, magnetosphere, and ionosphere. Space plasma interacts with spacecraft surfaces through electrostatic charging and current collection, potentially causing electrical discharges that damage electronics. Understanding plasma behavior is critical for designing spacecraft electrical systems, predicting radio communication blackouts during re-entry, and interpreting scientific measurements from space missions.}{catAE}

\dictentry{Space Vacuum}{The extremely low-pressure environment of outer space, with pressures typically below 10 to the minus 12 torr, where the absence of atmospheric molecules creates unique engineering challenges for spacecraft design. The vacuum causes materials to outgas, lubricants to evaporate, and heat transfer to occur only through radiation rather than convection. Spacecraft must be designed to operate in vacuum conditions with appropriately selected materials, sealed fluid systems, and thermal management approaches suited to the lack of convective cooling.}{catAE}

\dictentry{Space Weather}{Environmental conditions in near-Earth space driven by solar activity, including solar flares, coronal mass ejections, and geomagnetic storms that can disrupt satellite operations, communications, and navigation systems. Space weather effects include increased radiation exposure for crew, single-event upsets in electronics, atmospheric drag increases that alter satellite orbits, and ionospheric disturbances that degrade GPS accuracy. Space weather monitoring and forecasting are essential for protecting both crewed and uncrewed space assets.}{catAE}






\dictentry{Spin Recovery}{A set of established flight control inputs designed to exit an autorotative spin, an aggravated stall condition in which one wing is more deeply stalled than the other, causing the aircraft to rotate about a near-vertical axis while descending. Recovery typically involves applying full opposite rudder, moving the control column forward to break the stall, and then neutralizing controls as rotation stops to execute a pull-out from the resulting dive. Spin recovery techniques are specific to each aircraft type and must be practiced in the actual aircraft during flight testing.}{catAE}

\dictentry{Split Flap}{A trailing edge high-lift device that hinges downward from the lower surface of the wing without sliding aft, increasing camber and lift coefficient during takeoff and landing. The split flap creates a large separated wake behind the lower surface that augments lift but generates more drag than slotted or Fowler flap designs. Despite the higher drag penalty, split flaps are mechanically simple and lightweight, making them suitable for light aircraft and vintage designs where low-speed performance improvements are needed without complex mechanisms.}{catAE}

\dictentry{Staging}{The practice of discarding empty propellant tanks, engines, and structural elements during a rocket's ascent to reduce the remaining vehicle mass and improve overall performance. Each discarded stage eliminates dead weight that would otherwise require additional propellant to accelerate, allowing the remaining stages to achieve higher velocities. Staging can be performed in parallel, where multiple stages fire simultaneously, or in series, where each stage ignites after the previous one is jettisoned.}{catAE}

\dictentry{Stall Recovery}{A defined set of flight control inputs executed to restore airflow over the wings and regain controlled flight when the aircraft enters an aerodynamic stall, a condition where the wing's angle of attack exceeds the critical limit and lift is dramatically reduced. Recovery typically involves reducing the angle of attack by moving the control column forward while applying full power, then carefully leveling the wings as airspeed and lift are restored. Stall recovery techniques vary by aircraft type and must be practiced regularly by pilots to build muscle memory for this critical emergency maneuver.}{catAE}

\dictentry{Stall Warning System}{An avionics system that detects an impending aerodynamic stall by monitoring the wing's angle of attack or aerodynamic buffet characteristics and alerts the pilot through visual, aural, and tactile warnings before the stall fully develops. The system typically activates a stick shaker that vibrates the control column at a predefined angle of attack margin ahead of the actual stall. More advanced systems include stick pushers that automatically apply nose-down elevator to prevent the aircraft from exceeding the critical angle of attack.}{catAE}



\dictentry{State Estimation}{The mathematical process of inferring the current state of a dynamic system, such as an aircraft or spacecraft, by combining predictions from a dynamic model with measurements from multiple sensors that may be noisy, incomplete, or subject to time delays. State estimation algorithms such as Kalman filters and particle filters optimally weight the relative uncertainties of model predictions and sensor observations to produce the most probable estimate of position, velocity, attitude, and other state variables. In aerospace, accurate state estimation is fundamental to navigation, guidance, control, and fault detection across all vehicle types and mission phases.}{catAE}

\dictentry{Stick Shaker}{A mechanical device mounted on or near the pilot's control column that rapidly vibrates the stick to provide a tactile and audible warning of an approaching aerodynamic stall before the stall actually occurs. The stick shaker activates at a predefined angle of attack margin, typically several degrees below the critical stall angle, giving the pilot time to take corrective action. It serves as the primary stall warning device in transport category aircraft and is required by certification regulations to alert pilots during conditions where visual or aural cues alone might be missed.}{catAE}

\dictentry{Stringer}{A thin longitudinal structural member attached to the inner surface of the aircraft skin that reinforces the skin against buckling under compressive and shear loads while distributing aerodynamic forces to the transverse frames and ribs. Stringers are typically extruded from aluminum alloy or formed from composite material and run spanwise along the fuselage or wing skin. In semi-monocoque construction, stringers work in conjunction with the skin, frames, and bulkheads to carry flight loads, forming an efficient load-bearing structure that is both lightweight and damage-tolerant.}{catAE}

\dictentry{Structural Test}{A program of physical tests performed on aircraft components, subassemblies, and full-scale airframes to verify that the structure meets all design strength, stiffness, and durability requirements before certification. Structural tests include static proof loading to determine ultimate strength, fatigue testing to simulate a lifetime of cyclic loads, and damage tolerance testing to verify that the structure can sustain cracks without catastrophic failure. Test results are compared against analytical predictions and form a critical part of the certification data package submitted to aviation authorities for type approval.}{catAE}





\dictentry{Superalloy}{A high-performance metallic alloy specifically engineered to retain exceptional mechanical strength, creep resistance, and oxidation resistance at elevated temperatures exceeding 600 degrees Celsius where conventional alloys would rapidly degrade. Nickel-based superalloys are the most widely used in aerospace, forming the turbine blades, vanes, and combustion liners of jet engines where they endure extreme thermal and centrifugal stresses. The development of single-crystal superalloy casting techniques has enabled turbine blades to operate at temperatures approaching the melting point of the base metal through advanced internal cooling channel designs.}{catAE}

\dictentry{Supplemental Type Certificate}{An official approval issued by an aviation authority that authorizes a major modification or repair to a previously type-certificated aircraft design, without requiring a completely new type certification process. STCs are issued when the proposed modification meets all applicable airworthiness requirements and does not adversely affect the original design's compliance with certification standards. The STC holder assumes responsibility for the modification's design and must provide all necessary documentation, instructions, and continued airworthiness data to support the modification throughout the aircraft's operational life.}{catAE}

\dictentry{Sustainable Aviation Fuel}{A aviation fuel produced from non-petroleum feedstocks including waste oils, agricultural residues, municipal solid waste, or synthesized from captured carbon dioxide and green hydrogen, designed to reduce the lifecycle carbon emissions of flight. SAF can be used as a drop-in replacement or blended with conventional jet fuel without requiring modifications to aircraft engines or fuel distribution infrastructure. While current production volumes remain a small fraction of total aviation fuel demand, SAF represents one of the most promising near-term strategies for decarbonizing commercial aviation.}{catAE}

\dictentry{Swashplate}{A mechanical assembly in a helicopter rotor system that translates the pilot's non-rotating control inputs from the collective and cyclic sticks into rotating pitch commands for each main rotor blade through a system of pitch links and pushrods. The swashplate consists of two concentric plates, one rotating with the rotor and one stationary, connected by a bearing that allows relative rotation while transmitting tilt and axial displacements. It is the critical interface between the pilot's cockpit controls and the rotor blades, enabling both collective pitch changes for vertical flight and cyclic pitch variations for directional control.}{catAE}

\dictentry{System Display}{An electronic cockpit display that presents real-time information about the status and operation of aircraft systems including hydraulic, electrical, fuel, pneumatic, and environmental control systems in synoptic format. System displays consolidate data from numerous sensors and switches into intuitive graphical representations that allow pilots to quickly assess normal system operation and identify abnormal conditions. These displays are typically managed by the ECAM or EICAS systems and automatically present relevant system pages when a fault is detected, supporting rapid crew decision-making during abnormal situations.}{catAE}

\dictentry{Tail Rotor}{A small rotor mounted vertically on the tail boom of a conventional helicopter that produces a sideways thrust to counteract the torque reaction produced by the main rotor, preventing the fuselage from spinning in the opposite direction of the main rotor. The tail rotor also provides directional yaw control through anti-torque pedal inputs from the pilot, allowing the helicopter to turn and maintain heading during hover and low-speed flight. Tail rotor designs vary from fully articulated two-bladed systems to enclosed fenestron configurations, each offering different trade-offs in noise, efficiency, and safety.}{catAE}

\dictentry{Tailwheel}{A small steerable wheel mounted at the extreme aft end of the fuselage on conventional landing gear aircraft, providing ground support and directional steering during taxi, takeoff roll, and landing rollout. Tailwheel-equipped aircraft, also known as taildraggers, have a center of gravity located behind the main gear, making them more challenging to handle on the ground due to the tendency to ground loop. Despite the greater pilot skill required, tailwheel landing gear offers advantages in weight, simplicity, and ground clearance for propeller-driven aircraft operating from unimproved strips.}{catAE}

\dictentry{TCAS}{Traffic Alert and Collision Avoidance System, an onboard avionics system that uses transponder signals from nearby aircraft to detect potential mid-air collisions and provides pilots with traffic advisories and resolution advisories to maintain safe separation. TCAS operates independently of air traffic control, interrogating the Mode C and Mode S transponders of nearby aircraft to determine their range, bearing, altitude, and closing rate. When a collision threat is detected, TCAS issues visual and aural alerts and commands specific vertical maneuvers that are coordinated between conflicting aircraft to ensure mutual avoidance.}{catAE}



\dictentry{Throttle}{The primary engine power control lever in the cockpit that the pilot moves to regulate the amount of fuel delivered to the engine, thereby controlling thrust output. Throttle levers are typically located in a central quadrant between the pilots and are connected to the engine fuel control unit through mechanical linkages, cables, or electronic signals in full-authority digital engine control systems. Modern throttle levers incorporate detents for specific power settings such as idle, climb, and takeoff power, and may include autothrottle servos that automatically adjust thrust during automated flight phases.}{catAE}



\dictentry{Tiltrotor}{A vertical takeoff and landing aircraft that uses engine-driven rotors mounted on tilting nacelles at the wing tips, allowing the rotors to be oriented vertically for helicopter-like hover and takeoff, then tilted forward for high-speed fixed-wing cruise flight. The tiltrotor combines the versatility of a helicopter with the speed, range, and efficiency of a conventional airplane, making it suitable for military transport and urban air mobility applications. The Bell Boeing V-22 Osprey is the most prominent operational tiltrotor, demonstrating the capability to cruise at speeds exceeding 500 kilometers per hour while retaining vertical flight capability.}{catAE}

\dictentry{Tip Speed}{The linear velocity of the helicopter rotor blade tip as it rotates, determined by the product of the rotational speed and the blade radius, typically ranging from 200 to 250 meters per second for conventional helicopters. Tip speed is a critical design parameter because as it approaches the speed of sound, compressibility effects including shock waves and drag divergence cause significant performance degradation and noise generation. Managing tip speed involves balancing the need for sufficient rotor thrust against the penalties of high-speed aerodynamic effects, and is particularly important during fast forward flight where the advancing blade tip may approach supersonic speeds.}{catAE}

\dictentry{Titanium Alloy}{A high-strength, corrosion-resistant metallic alloy with titanium as the primary element, valued in aerospace for its excellent strength-to-weight ratio at both ambient and elevated temperatures up to approximately 600 degrees Celsius. Titanium alloys are used for critical structural components including landing gear, engine fan blades, compressor discs, and fasteners where the combination of strength, temperature resistance, and corrosion immunity justifies the higher material and machining costs. The most common aerospace titanium alloy, Ti-6Al-4V, accounts for roughly half of all titanium usage in aircraft and spacecraft structures.}{catAE}

\dictentry{Trailing Edge Device}{An aerodynamic mechanism mounted on the trailing edge of a wing that modifies the wing's camber, chord, or surface area to increase lift for low-speed flight during takeoff and landing. Trailing edge devices include plain flaps, split flaps, slotted flaps, Fowler flaps, and spoilers, each offering different combinations of lift increment, drag, and mechanical complexity. When retracted during cruise, these devices fit flush with the wing profile to maintain an efficient aerodynamic shape, and they are extended selectively during low-speed flight phases to reduce approach speeds and shorten runway requirements.}{catAE}

\dictentry{Transfer Orbit}{An intermediate orbital path used to move a spacecraft between two different orbits, most commonly an elliptical trajectory connecting a lower parking orbit to a higher target orbit through one or more propulsive burns. The Hohmann transfer is the most fuel-efficient two-burn transfer between coplanar circular orbits, while bi-elliptic transfers can be more efficient for large orbit changes. Transfer orbits are fundamental to mission planning for geostationary orbit insertion, interplanetary trajectories, and orbital rendezvous operations.}{catAE}

\dictentry{Transformer}{An electrical device that transfers alternating current energy between two or more circuits through electromagnetic induction, changing voltage and current levels while maintaining the same frequency. In aircraft electrical systems, transformers step up or step down voltages to match the requirements of different loads, isolate circuits for safety, and provide multiple voltage outputs from a single source. Modern aerospace transformers use lightweight core materials and advanced winding techniques to minimize weight while meeting stringent reliability requirements for continuous operation across wide temperature and vibration ranges.}{catAE}

\dictentry{Translational Lift}{The increase in helicopter rotor efficiency and lift-to-power ratio that occurs during forward flight as the rotor moves through undisturbed air rather than operating in its own recirculating downwash. At hover the rotor must overcome its own induced velocity field, but as forward speed increases, each blade section encounters progressively cleaner air, reducing induced drag and improving aerodynamic performance. This effect is most pronounced at low to moderate forward speeds and diminishes at higher speeds where parasite drag becomes the dominant power requirement.}{catAE}



\dictentry{Tricycle Gear}{A landing gear configuration consisting of a single nose wheel or nose gear assembly positioned forward of the aircraft's center of gravity and two or more main gear units located aft near the center of gravity, providing stable ground handling and preventing ground looping tendencies. Tricycle gear is the dominant landing gear arrangement for modern aircraft because it allows the fuselage to remain level during taxi, improving pilot visibility and passenger comfort. The nose gear is typically steerable through the rudder pedals or a tiller wheel, providing precise directional control during ground operations.}{catAE}

\dictentry{Trim Tab}{A small hinged surface attached to the trailing edge of a primary control surface that the pilot adjusts to relieve sustained control forces, allowing the aircraft to maintain a desired attitude without continuous manual input. Trim tabs can be manually adjusted from the cockpit using a trim wheel or electric switch, and they work by creating a small aerodynamic force that counteracts the aerodynamic imbalance on the main control surface. Proper trim reduces pilot fatigue, improves fuel efficiency by maintaining optimal attitudes, and ensures the aircraft can be flown hands-off during steady-state flight conditions.}{catAE}

\dictentry{True Airspeed}{The actual speed of the aircraft relative to the surrounding air mass, corrected for the effects of air density variations with altitude and temperature, which differs from both indicated and calibrated airspeed. True airspeed increases with altitude for a given indicated airspeed because the thinner air produces less dynamic pressure on the airspeed indicator. Pilots use true airspeed for fuel planning, time and distance calculations, and wind correction computations, and it is typically displayed on modern air data computers and flight management systems.}{catAE}

\dictentry{True Heading}{The direction in which the aircraft's longitudinal axis is pointed, measured as an angle clockwise from true north, which differs from magnetic heading by the local magnetic variation or declination. True heading is essential for accurate navigation over long distances, as it relates directly to geographic coordinates and great circle route calculations. Pilots convert between true and magnetic heading using the local variation value, which changes with geographic location and must be updated as the aircraft traverses different regions.}{catAE}





\dictentry{Turboprop}{A gas turbine engine in which the turbine extracts most of the energy from the hot gas stream to drive a propeller through a reduction gearbox, with only a small amount of residual thrust produced by the exhaust nozzle. Turboprops offer excellent fuel efficiency at low to moderate airspeeds and altitudes, making them ideal for regional airliners, military transports, and utility aircraft that operate from short runways. The reduction gearbox is a critical and complex component that steps down the high turbine RPM to an optimal propeller speed, typically in the range of 1,200 to 2,000 RPM.}{catAE}

\dictentry{Turboshaft}{A gas turbine engine optimized to deliver shaft power rather than jet thrust, in which nearly all of the turbine energy is extracted to drive an output shaft connected to a helicopter transmission, generator, or industrial load. Turboshaft engines are the primary power plant for most modern helicopters, offering high power-to-weight ratios, reliable operation, and the ability to burn readily available jet fuel. The engine's power output is controlled by a fuel control unit that adjusts fuel flow to maintain a constant output shaft speed regardless of the load demand from the rotor system.}{catAE}

\dictentry{Type Certificate}{An official document issued by an aviation authority such as the FAA or EASA certifying that an aircraft design meets all applicable airworthiness, noise, and emissions standards and is safe for operation within its approved limitations. The type certificate is awarded after an exhaustive certification process that includes ground testing, flight testing, and review of all design documentation, manufacturing processes, and maintenance procedures. Once issued, the type certificate defines the aircraft's approved configuration, operating limitations, and continued airworthiness requirements that must be maintained throughout the product's operational life.}{catAE}


\dictentry{Van Allen Belts}{Two distinct toroidal zones of energetic charged particles trapped by the Earth's magnetic field, consisting of an inner belt primarily of high-energy protons at altitudes of roughly 1,000 to 6,000 kilometers and an outer belt of high-energy electrons extending from about 13,000 to 60,000 kilometers. The Van Allen belts pose significant radiation hazards to spacecraft electronics and crew, requiring radiation-hardened components and shielding for missions that transit or operate within these regions. Spacecraft in low Earth orbit typically pass through only the weaker fringe of the inner belt, while missions to geostationary orbit or beyond must navigate through both belts during ascent and transit.}{catAE}

\dictentry{Ventilation}{The controlled circulation and replacement of air within the aircraft cabin and cockpit to maintain air quality, remove contaminants, and ensure occupant comfort and safety during all phases of flight. The ventilation system works in conjunction with the environmental control system to supply fresh air, remove carbon dioxide and odors, and maintain appropriate humidity levels throughout the passenger and crew compartments. Adequate ventilation is critical not only for comfort but also for preventing the buildup of harmful fumes from cabin materials, food preparation, and potential engine bleed air contamination events.}{catAE}

\dictentry{Vertical Stabilizer}{The fixed vertical tail surface mounted at the rear of the aircraft that provides directional stability by generating a restoring side force and yawing moment when the aircraft sideslips or weathers off its heading. The vertical stabilizer serves as the mounting structure for the movable rudder control surface, which the pilot uses to control yaw and coordinate turns. Its size and placement are determined by the aircraft's directional stability requirements and must provide adequate control authority during engine-out conditions on multi-engine aircraft.}{catAE}



\dictentry{Voltage Regulator}{An electrical control device that maintains a constant output voltage level regardless of variations in input voltage, load current, or temperature, ensuring that sensitive avionics and electrical equipment receive stable and clean power. In aircraft systems, voltage regulators are integrated into alternators and generators to control field current and maintain the electrical bus voltage within narrow limits specified by aerospace standards. Modern solid-state voltage regulators offer fast response times, high reliability, and precise regulation, protecting downstream equipment from voltage spikes and sags that occur during transient load changes.}{catAE}

\dictentry{Vortex Ring State}{A dangerous helicopter aerodynamic condition, also called settling with power, in which the rotor becomes engulfed in its own turbulent downwash during a steep descent at low airspeed, causing a severe loss of lift and an uncontrollable rate of descent. The helicopter effectively settles into a recirculating vortex ring that surrounds the rotor disc, and applying more power only intensifies the downwash and worsens the condition. Recovery requires the pilot to tilting the cyclic forward to accelerate out of the vortex ring and transition into forward flight where clean airflow restores normal rotor performance.}{catAE}

\dictentry{Waypoint}{A predetermined geographical position defined by specific latitude and longitude coordinates that serves as a reference point along a planned flight route, used for navigation, air traffic control separation, and flight management system programming. Waypoints are typically established at intersections of airways, holding patterns, approach fixes, and points of interest, and may be named using five-letter alphanumeric identifiers standardized by ICAO. Modern flight management systems sequence through waypoints automatically, computing optimal course, speed, and altitude profiles between successive points to achieve efficient route compliance.}{catAE}

\dictentry{Wing Spar}{The principal spanwise structural member of a wing that carries the primary bending and shear loads generated by aerodynamic forces and the aircraft's weight, extending from the wing root at the fuselage to the wing tip. Wing spars are typically constructed as box beams or I-beams from high-strength aluminum alloy or composite materials, with the spar caps carrying bending loads and the spar webs resisting shear forces. The spar also serves as the primary attachment structure for engine mounts, landing gear, control surfaces, fuel tanks, and other wing-mounted equipment.}{catAE}

\dictentry{Yaw Damper}{An automatic flight control system component that detects and suppresses Dutch roll oscillations, a coupled lateral-directional oscillation in which the aircraft alternately yaws and rolls in a repeating cycle that is uncomfortable for passengers and potentially hazardous if uncontrolled. The yaw damper uses a yaw rate gyroscope to sense the oscillatory yaw motion and commands small, rapid rudder inputs through the autopilot servos to dampen the oscillation before it becomes noticeable. In transport category aircraft, the yaw damper operates continuously during flight and is required by certification regulations for aircraft with certain directional stability characteristics.}{catAE}

\dictentry{Yoke}{The primary flight control input device in the cockpit, resembling a steering wheel or W-shaped handle mounted on a column, that the pilot manipulates to command pitch and roll of the aircraft. Pulling the yoke aft raises the elevator and pitches the nose up, while pushing forward lowers the elevator for nose-down pitch, and rotating the yoke left or right deflects the ailerons to roll the aircraft in the corresponding direction. The yoke may also incorporate switches for trim, radio communication, and autopilot disconnect, and in some designs it is connected to the control surfaces through a combination of mechanical cables and hydraulic boost systems.}{catAE}

